\chapter{Complex Analysis}\label{complex-analysis}
\chapterstatus{complete}

\section{Introduction}\label{complex-analysis:section-introduction}

A function of one complex variable that is differentiable once is analytic, and is determined by its values on any set with a limit
point. This rigidity, which has no analogue in real analysis, is the source of the special features of complex geometry: on a compact
complex manifold there are few holomorphic functions, and the interplay between holomorphic and topological invariants leads to Hodge
theory. This chapter proves Cauchy's theorem in homotopy form (Theorem~\ref{complex-analysis:theorem-cauchy-homotopy}), the Cauchy
integral formula and its consequences, the residue theorem, and the Riemann mapping theorem
(Theorem~\ref{complex-analysis:theorem-riemann-mapping}). The last section studies the Riemann sphere $\hat{\CC} = \CC \cup \{\infty\}$: stereographic
projection, the rational functions as the meromorphic functions on $\hat{\CC}$ (Theorem~\ref{complex-analysis:theorem-rational}), and the M\"obius transformations, with the
cross-ratio and their classification into parabolic, elliptic, hyperbolic and loxodromic ones (Theorem~\ref{complex-analysis:theorem-mobius-classification}). References: \cite{ahlfors-complex}, \cite{stein-shakarchi-complex},
\cite{conway-complex} and \cite[Chapters 10--16]{rudin-rca}.

\noindent\textbf{Prerequisites.} Chapter~\ref{real-analysis} (Real Analysis) and Chapter~\ref{homotopy} (Homotopy and Covering Spaces).

\noindent\textbf{On the way to the Hodge conjecture.} Needed: holomorphic functions, Cauchy's theorem and the Cauchy integral formula, and Laurent series and residues. The Riemann mapping theorem
and the Riemann sphere (Sections~\ref{complex-analysis:section-riemann-mapping}, \ref{complex-analysis:section-riemann-sphere}) can be postponed.

We write $D(a, r) = \{z : |z - a| < r\}$, $\overline{D}(a, r)$ for its closure, and $\mathbb{D} = D(0, 1)$. All domains ($=$ nonempty connected
open subsets of $\CC$) are denoted $U$, $V$.

\section{Holomorphic functions}\label{complex-analysis:section-holomorphic}

\begin{definition}\label{complex-analysis:definition-holomorphic}
Let $U \subseteq \CC$ be open. A function $f \colon U \to \CC$ is \emph{holomorphic} if the limit $f'(a) = \lim_{h \to 0}(f(a + h) - f(a))/h$, $h \in \CC^\times$, exists for every $a \in U$. An
\emph{entire} function is one holomorphic on $\CC$.
\end{definition}

The definition looks exactly like the one for real functions, but $h$ now tends to $0$ through all complex numbers, from every direction at once, and the difference quotient must have the same limit
along all of them. This single requirement turns out to be extraordinarily strong: it implies that $f$ has derivatives of all orders and is given by its Taylor series
(Theorem~\ref{complex-analysis:theorem-analytic} and Corollary~\ref{complex-analysis:corollary-smooth-harmonic}). Before seeing why, let us compute a few derivatives directly from the definition.

\begin{example}\label{complex-analysis:example-derivative-definition}
Let $a \in \CC$ and $h \in \CC^\times$ small.
\begin{enumerate}
\item $f(z) = z^n$, $n \geq 1$. By the binomial theorem $(a + h)^n - a^n = na^{n-1}h + \sum_{k = 2}^n\binom nka^{n-k}h^k$, so the difference quotient is
$na^{n-1} + \sum_{k = 2}^n\binom nka^{n-k}h^{k-1} \to na^{n-1}$. So $(z^n)^{\prime} = nz^{n-1}$, as over $\RR$.
\item $f(z) = 1/z$ on $\CC^\times$. For $a \neq 0$ and $|h| < |a|$, $\frac1h\bigl(\frac1{a + h} - \frac1a\bigr) = \frac{-1}{a(a + h)} \to -\frac1{a^2}$.
\item $f(z) = \bar z$. The difference quotient is $\bar h/h$, which equals $1$ for real $h$ and $-1$ for purely imaginary $h$. So the limit exists at no point: $\bar z$ is a perfectly smooth map of the
plane, an isometry even, but it is nowhere complex differentiable. The same happens for $\operatorname{Re}z$, whose quotient $\operatorname{Re}h/h$ is $1$ on the real axis and $0$ on the imaginary
axis.
\item $f(z) = |z|^2 = z\bar z$. Here $|a + h|^2 - |a|^2 = \bar ah + a\bar h + |h|^2$, so the quotient is $\bar a + a\,\bar h/h + \bar h$. For $a = 0$ it is $\bar h \to 0$; for $a \neq 0$ the middle term
$a\,\bar h/h$ has no limit, by (3). So $|z|^2$ is complex differentiable at $0$ and nowhere else. A function can be complex differentiable at a single point; ``holomorphic'' asks for it at every point
of an open set, and all the theory below depends on this.
\end{enumerate}
\end{example}

\begin{proposition}[Cauchy--Riemann]\label{complex-analysis:proposition-cauchy-riemann}
Identify $\CC$ with $\RR^2$ and write $f = u + iv$. Then $f$ is holomorphic at $a$ if and only if $f$ is differentiable at $a$ as a map of real variables
(Definition~\ref{real-analysis:definition-differentiable}) and its derivative is complex linear, which holds if and only if
\[
\partial_xu = \partial_yv, \qquad \partial_yu = -\partial_xv \quad \text{at } a.
\]
Then $f'(a) = \partial_xu + i\partial_xv$. Sums, products, quotients and composites of holomorphic functions are holomorphic, with the usual rules.
\end{proposition}

\begin{proof}
Complex differentiability at $a$ says $f(a + h) = f(a) + ch + o(|h|)$ with $c = f'(a) \in \CC$, which is real differentiability with $Df(a)$ the real linear map $h \mapsto ch$, that is, with matrix
$\left(\begin{smallmatrix} \alpha & -\beta \\ \beta & \alpha \end{smallmatrix}\right)$ for $c = \alpha + i\beta$. Comparing with the Jacobian matrix
$\left(\begin{smallmatrix} \partial_xu & \partial_yu \\ \partial_xv & \partial_yv \end{smallmatrix}\right)$ gives the equations. The rules follow as in
Proposition~\ref{real-analysis:proposition-derivative-rules}, whose proofs are valid over $\CC$.
\end{proof}

\begin{example}\label{complex-analysis:example-holomorphic}
Polynomials are entire; $1/z$ is holomorphic on $\CC^\times$; power series are holomorphic in their disc of convergence, with the termwise derivative
(Theorem~\ref{real-analysis:theorem-power-series-derivative}); $\exp$ is entire with $\exp' = \exp$. The function $\bar z$ is nowhere holomorphic, since $\partial_xu = 1 \neq -1 = \partial_yv$.
\end{example}

The Cauchy--Riemann equations have a geometric meaning. A real linear map $\RR^2 \to \RR^2$ is complex linear exactly when its matrix is
$\left(\begin{smallmatrix} \alpha & -\beta \\ \beta & \alpha \end{smallmatrix}\right)$, that is, when it is multiplication by $c = \alpha + i\beta = |c|e^{i\varphi}$: a rotation by $\varphi$ followed by
a dilation by $|c|$. So near a point where $f^{\prime}(a) \neq 0$, a holomorphic map looks, to first order, like a rotation and a dilation, which preserve angles and orientation. Maps with this property are
called \emph{conformal}, and the converse also holds.

\begin{proposition}\label{complex-analysis:proposition-conformal}
\begin{enumerate}
\item Let $f$ be holomorphic near $a$ with $f^{\prime}(a) \neq 0$, and let $\gamma_1, \gamma_2$ be $C^1$ curves with $\gamma_1(0) = \gamma_2(0) = a$ and nonzero tangent vectors $\gamma_1^{\prime}(0), \gamma_2^{\prime}(0)$.
Then $(f \circ \gamma_j)^{\prime}(0) = f^{\prime}(a)\gamma_j^{\prime}(0)$. Hence the image curves have tangent vectors rotated by the same angle $\arg f^{\prime}(a)$ and stretched by the same factor
$|f^{\prime}(a)|$, and the oriented angle between them equals that between $\gamma_1$ and $\gamma_2$.
\item Conversely, let $f$ be real differentiable at $a$ with $A = Df(a)$ invertible, and suppose that $A$ preserves oriented angles, that is, $Aw/Av$ and $w/v$ have the same argument for all nonzero
$v, w \in \CC$. Then $f$ is complex differentiable at $a$.
\end{enumerate}
\end{proposition}

\begin{proof}
(1) By the chain rule (Proposition~\ref{real-analysis:proposition-derivative-rules}) and Proposition~\ref{complex-analysis:proposition-cauchy-riemann},
$(f \circ \gamma_j)^{\prime}(0) = Df(a)\gamma_j^{\prime}(0) = f^{\prime}(a)\gamma_j^{\prime}(0)$. The quotient of the two image tangent vectors is $\gamma_2^{\prime}(0)/\gamma_1^{\prime}(0)$, the same as for the original
ones, and the oriented angle from one vector to another is the argument of this quotient.

(2) Let $p = A(1) \neq 0$ and $q = A(i)$. The oriented angle from $1$ to $i$ is $\frac\pi2$, so $q/p$ has argument $\frac\pi2$: $q = itp$ for some $t > 0$. The oriented angle from $1$ to $1 + i$ is $\frac\pi4$, so
$A(1 + i)/A(1) = (p + q)/p = 1 + it$ has argument $\frac\pi4$, which forces $t = 1$. Hence $A(x + iy) = xp + yip = (x + iy)p$: $A$ is multiplication by $p$, so it is complex linear, and $f$ is
complex differentiable at $a$ with $f^{\prime}(a) = p$ by Proposition~\ref{complex-analysis:proposition-cauchy-riemann}.
\end{proof}

\begin{example}\label{complex-analysis:example-conformal}
\begin{enumerate}
\item \emph{Critical points.} For $f(z) = z^2$ and $a = 0$, where $f^{\prime}(0) = 0$, angles are not preserved but doubled: the ray $\{re^{i\theta} : r > 0\}$ is mapped to the ray $\{r^2e^{2i\theta}\}$. In
general, near a zero of order $m$ of $f^{\prime}$ the map multiplies angles by $m + 1$; this follows from the local form of holomorphic maps (Theorem~\ref{complex-analysis:theorem-open-mapping}).
\item \emph{The exponential.} $e^{x + iy} = e^xe^{iy}$ maps the vertical line $x = x_0$ onto the circle of radius $e^{x_0}$, and the horizontal line $y = y_0$ onto the ray of argument $y_0$. The lines of the
grid meet at right angles, and so do their images, circles and rays from the origin, as Proposition~\ref{complex-analysis:proposition-conformal} predicts, since $\exp^{\prime} = \exp$ never vanishes.
\item \emph{Reflection.} $z \mapsto \bar z$ preserves the size of angles but reverses their orientation, and it is not holomorphic
(Example~\ref{complex-analysis:example-derivative-definition}(3)). The argument in (2) adapts: an invertible real linear map that preserves unoriented angles is either complex linear or of the form
$z \mapsto p\bar z$, so the maps that preserve unoriented angles, at points where the derivative is invertible, are the holomorphic and the \emph{antiholomorphic} ones (those $f$ with $\bar f$ holomorphic).
\end{enumerate}
\end{example}

\begin{definition}\label{complex-analysis:definition-path-integral}
A \emph{path} is a piecewise continuously differentiable map $\gamma \colon [a, b] \to \CC$; it is \emph{closed} if $\gamma(a) = \gamma(b)$. Its image is $\gamma^*$ and its length is
$L(\gamma) = \int_a^b|\gamma'(t)|\,dt$. For continuous $f$ on $\gamma^*$,
\[
\int_\gamma f(z)\,dz = \int_a^bf(\gamma(t))\gamma'(t)\,dt.
\]
\end{definition}

\begin{lemma}\label{complex-analysis:lemma-path-integral}
The integral is unchanged by increasing $C^1$ reparametrisations, changes sign under reversal, and $|\int_\gamma f| \leq L(\gamma)\sup_{\gamma^*}|f|$. If $f = F'$ for a holomorphic $F$ on an open set
containing $\gamma^*$, then $\int_\gamma f = F(\gamma(b)) - F(\gamma(a))$; in particular $\int_\gamma f = 0$ for closed $\gamma$. For the circle $\gamma(t) = a + re^{2\pi it}$, $t \in [0, 1]$, one has
$\int_\gamma(z - a)^n\,dz = 0$ for $n \neq -1$ and $= 2\pi i$ for $n = -1$.
\end{lemma}

\begin{proof}
Reparametrisation and reversal follow from the substitution rule (Corollary~\ref{real-analysis:corollary-integration-rules}), the estimate from
Proposition~\ref{real-analysis:proposition-riemann}. If $f = F'$, then $(F \circ \gamma)' = f(\gamma)\gamma'$ on each $C^1$ piece, so the integral telescopes by the fundamental theorem of calculus. For the
circle and $n \neq -1$, $(z - a)^n$ has the primitive $(z - a)^{n+1}/(n + 1)$; for $n = -1$ the integral is $\int_0^12\pi i\,dt$.
\end{proof}

In the definition, $f(\gamma(t))\gamma^{\prime}(t)$ is a complex-valued function of the real variable $t$, and its integral is taken separately on the real and the imaginary part
(Definition~\ref{real-analysis:definition-riemann}): $\int_a^b(g + ih) = \int_a^bg + i\int_a^bh$. So a complex line integral is a pair of real line integrals, and the formal rule
$f\,dz = (u + iv)(dx + i\,dy)$ says which ones.

\begin{proposition}\label{complex-analysis:proposition-path-integral-real}
Let $\gamma = x + iy$ be a path and $f = u + iv$ continuous on $\gamma^*$, with $x, y, u, v$ real. Then
\[
\int_\gamma f(z)\,dz = \int_\gamma(u\,dx - v\,dy) + i\int_\gamma(v\,dx + u\,dy),
\]
where the two integrals on the right are real line integrals of the second kind (Definition~\ref{real-analysis:definition-second-kind}) of the vector fields $(u, -v)$ and $(v, u)$ in $\RR^2 = \CC$. Moreover
$\bigl|\int_\gamma f(z)\,dz\bigr| \leq \int_\gamma|f|\,ds$, where $\int_\gamma|f|\,ds = \int_a^b|f(\gamma(t))|\,|\gamma^{\prime}(t)|\,dt$ is a line integral of the first kind
(Definition~\ref{real-analysis:definition-first-kind}), also written $\int_\gamma|f|\,|dz|$.
\end{proposition}

\begin{proof}
On each $C^1$ piece, $f(\gamma(t))\gamma^{\prime}(t) = (u + iv)(x^{\prime} + iy^{\prime}) = (ux^{\prime} - vy^{\prime}) + i(vx^{\prime} + uy^{\prime})$, with $u$ and $v$ evaluated at $\gamma(t)$. The real part is
$\langle(u, -v), (x^{\prime}, y^{\prime})\rangle$ and the imaginary part $\langle(v, u), (x^{\prime}, y^{\prime})\rangle$, and integrating them gives the two line integrals. The estimate is the norm inequality
$|\int g| \leq \int|g|$ of Proposition~\ref{real-analysis:proposition-riemann}(3), since $|f(\gamma(t))\gamma^{\prime}(t)| = |f(\gamma(t))|\,|\gamma^{\prime}(t)|$.
\end{proof}

For instance, for $f(z) = \bar z = x - iy$ the two real fields are $(x, y)$ and $(-y, x)$. The first is the gradient of $\frac12(x^2 + y^2)$, so its integral over a closed path vanishes; the second is the rotation field,
whose integral over the unit circle is $2\pi$ (Example~\ref{real-analysis:example-second-kind}(1)). This gives $\int_\gamma\bar z\,dz = 0 + 2\pi i$, as in Example~\ref{complex-analysis:example-path-integrals}(1) below.

\begin{remark}\label{complex-analysis:remark-cauchy-green}
The two real forms are closed exactly when $f$ satisfies the Cauchy--Riemann equations: for $u, v$ of class $C^1$, $d(u\,dx - v\,dy) = -(u_y + v_x)\,dx \wedge dy$ and $d(v\,dx + u\,dy) = (u_x - v_y)\,dx \wedge dy$
(Proposition~\ref{complex-analysis:proposition-cauchy-riemann}). So, if $f$ is holomorphic with continuous derivative and $\gamma$ bounds a region $D$ in its domain, Green's theorem (Chapter~\ref{riemannian-geometry})
gives $\int_\gamma f\,dz = \iint_D0 = 0$: this is Cauchy's original proof of his theorem. The proof below, through Goursat's theorem, does not assume that $f^{\prime}$ is continuous, and does not need Green's
theorem. In the language of Chapter~\ref{differential-forms}, $f\,dz$ is a complex-valued $1$-form, and $d(f\,dz) = \frac{\partial f}{\partial\bar z}\,d\bar z \wedge dz$, which vanishes for holomorphic $f$.
\end{remark}

\begin{example}\label{complex-analysis:example-path-integrals}
\begin{enumerate}
\item \emph{A function without primitive.} Let $\gamma(t) = e^{2\pi it}$, $t \in [0, 1]$, be the unit circle. On it $\bar z = e^{-2\pi it}$ and $\gamma^{\prime}(t) = 2\pi ie^{2\pi it}$, so
$\int_\gamma\bar z\,dz = \int_0^12\pi i\,dt = 2\pi i \neq 0$. By Lemma~\ref{complex-analysis:lemma-path-integral}, $\bar z$ is not the derivative of a holomorphic function on any open set containing the
circle.
\item \emph{Dependence on the path.} Integrate $\bar z$ from $0$ to $1 + i$ in two ways. Along the segment $\gamma(t) = t(1 + i)$: $\int_0^1t(1 - i)(1 + i)\,dt = \int_0^12t\,dt = 1$. Along the segment from
$0$ to $1$ followed by the segment from $1$ to $1 + i$, parametrised by $t$ and $1 + it$: $\int_0^1t\,dt + \int_0^1(1 - it)\,i\,dt = \frac12 + \bigl(i + \frac12\bigr) = 1 + i$. The results differ. For a
function with a primitive, such as $z^2 = (z^3/3)^{\prime}$, every path from $0$ to $1 + i$ gives $(1 + i)^3/3$.
\item \emph{A holomorphic function without primitive.} $1/z$ is holomorphic on $\CC^\times$, yet $\int_\gamma dz/z = 2\pi i$ over the unit circle (Lemma~\ref{complex-analysis:lemma-path-integral}); so $1/z$ has no
primitive on $\CC^\times$, that is, there is no holomorphic logarithm on $\CC^\times$. What goes wrong is the hole at $0$: Cauchy's theorem below says that integrals of holomorphic functions over closed
paths vanish when the path can be shrunk to a point inside the domain.
\end{enumerate}
\end{example}

\section{Cauchy's theorem}\label{complex-analysis:section-cauchy}

Example~\ref{complex-analysis:example-path-integrals} shows that integrals over closed paths need not vanish, either because the function is not holomorphic ($\bar z$) or because the domain has a hole
($1/z$). Cauchy's theorem says that these are the only obstructions. For a function $f = u + iv$ of class $C^1$ on a neighbourhood of a triangle $T$ there is a quick argument, using Green's formula
$\int_{\partial T}(P\,dx + Q\,dy) = \iint_T(\partial_xQ - \partial_yP)\,dx\,dy$ (proved in Chapter~\ref{differential-forms} in the language of forms) with $P = f$ and $Q = if$, since $dz = dx + i\,dy$:
\[
\int_{\partial T}f\,dz = \iint_T\bigl(i\,\partial_xf - \partial_yf\bigr)\,dx\,dy = 0,
\]
because the Cauchy--Riemann equations say exactly that $\partial_yf = i\,\partial_xf$ (Proposition~\ref{complex-analysis:proposition-cauchy-riemann}: $\partial_yf = \partial_yu + i\partial_yv = -\partial_xv + i\partial_xu$).
Goursat's theorem proves the same conclusion assuming only that $f^{\prime}$ exists, without its continuity, by subdividing the triangle and using the linear approximation of $f$ at one point. This is
what later allows us to prove, rather than assume, that holomorphic functions are smooth.

\begin{theorem}[Goursat]\label{complex-analysis:theorem-goursat}
Let $U$ be open, $p \in U$, and $f \colon U \to \CC$ continuous on $U$ and holomorphic on $U \setminus \{p\}$. Then $\int_{\partial T}f = 0$ for every closed triangle $T \subseteq U$, where $\partial T$ is its
boundary traversed once.
\end{theorem}

\begin{proof}
First assume $f$ is holomorphic on all of $U$. Let $I = \int_{\partial T}f$. Joining the midpoints of the sides divides $T$ into four triangles whose boundary integrals add up to $I$ (the inner edges cancel),
so one of them, $T_1$, has $|\int_{\partial T_1}f| \geq |I|/4$. Iterating gives nested triangles $T_n$ with $\operatorname{diam}T_n = 2^{-n}\operatorname{diam}T$, $L(\partial T_n) = 2^{-n}L(\partial T)$ and
$|\int_{\partial T_n}f| \geq 4^{-n}|I|$. Their intersection is a point $z_0$. Given $\varepsilon > 0$, holomorphy at $z_0$ gives $\delta$ with $|f(z) - f(z_0) - f'(z_0)(z - z_0)| \leq \varepsilon|z - z_0|$ for
$|z - z_0| < \delta$. The affine function $f(z_0) + f'(z_0)(z - z_0)$ has a primitive, so its integral over $\partial T_n$ vanishes, and for large $n$,
$|\int_{\partial T_n}f| \leq \varepsilon\operatorname{diam}(T_n)L(\partial T_n) = \varepsilon4^{-n}\operatorname{diam}(T)L(\partial T)$. Hence $|I| \leq \varepsilon\operatorname{diam}(T)L(\partial T)$ for all
$\varepsilon$.

If $f$ is only continuous at $p$: if $p \notin T$ we are done. If $p$ is a vertex of $T$, write $T$ as the union of a small triangle $T'$ with vertex $p$ and two triangles not containing $p$; then
$|\int_{\partial T}f| = |\int_{\partial T'}f| \leq L(\partial T')\sup_T|f| \to 0$ as $T'$ shrinks. For general $p \in T$, join $p$ to the vertices.
\end{proof}

\begin{corollary}\label{complex-analysis:corollary-primitive-convex}
Let $U$ be convex, $p \in U$, and $f$ continuous on $U$ and holomorphic on $U \setminus \{p\}$. Then $f$ has a primitive on $U$, and $\int_\gamma f = 0$ for every closed path $\gamma$ in $U$.
\end{corollary}

\begin{proof}
Fix $z_0 \in U$ and set $F(z) = \int_{[z_0, z]}f$, the integral over the segment. For $z, z + h \in U$, the triangle with vertices $z_0, z, z + h$ lies in $U$, so $F(z + h) - F(z) = \int_{[z, z+h]}f$ by
Theorem~\ref{complex-analysis:theorem-goursat}. Hence $|F(z + h) - F(z) - f(z)h| \leq |h|\sup_{[z, z+h]}|f - f(z)| = o(|h|)$ by continuity, so $F' = f$. Apply Lemma~\ref{complex-analysis:lemma-path-integral}.
\end{proof}

\begin{theorem}[Cauchy's theorem, homotopy form]\label{complex-analysis:theorem-cauchy-homotopy}
Let $f$ be holomorphic on an open $U \subseteq \CC$ and let $\gamma_0, \gamma_1$ be paths in $U$ that are homotopic in $U$, either as closed paths or with fixed endpoints. Then
$\int_{\gamma_0}f = \int_{\gamma_1}f$. In particular $\int_\gamma f = 0$ for every closed path that is null-homotopic in $U$, and every holomorphic function on a simply connected open set has a primitive.
\end{theorem}

\begin{proof}
Let $H \colon [0, 1]^2 \to U$ be the homotopy, $H(s, 0) = \gamma_0(s)$, $H(s, 1) = \gamma_1(s)$. Its image is compact, so there is $\rho > 0$ with $D(H(s, t), \rho) \subseteq U$ for all $(s, t)$, and by uniform
continuity there is $N$ with $|H(s, t) - H(s', t')| < \rho$ whenever $|s - s'|, |t - t'| \leq 1/N$. Put $z_{jk} = H(j/N, k/N)$ and let $D_{jk} = D(z_{jk}, \rho)$, so that $H$ maps the square
$[j/N, (j+1)/N] \times [k/N, (k+1)/N]$ into $D_{jk}$. Let $P_k$ be the closed polygonal path through $z_{0k}, z_{1k}, \ldots, z_{Nk}$.

Each segment $[z_{jk}, z_{j+1,k}]$ and the corresponding piece of $H(-, k/N)$ lie in the convex set $D_{jk}$, where $f$ has a primitive by Corollary~\ref{complex-analysis:corollary-primitive-convex}; so the
integrals of $f$ over them agree, and $\int_{\gamma_0}f = \int_{P_0}f$, $\int_{\gamma_1}f = \int_{P_N}f$ (for $\gamma_0$, $\gamma_1$ piecewise $C^1$; the intermediate paths are only continuous, and the
polygonal paths replace them). The closed polygon with vertices $z_{jk}, z_{j+1,k}, z_{j+1,k+1}, z_{j,k+1}$ lies in the convex set $D_{jk}$, so the integral of $f$ over it vanishes. Summing over $j$, the
contributions of the segments $[z_{jk}, z_{j,k+1}]$ cancel in pairs (for closed homotopies because $z_{0k} = z_{Nk}$; for fixed endpoints because those segments are constant), so
$\int_{P_k}f = \int_{P_{k+1}}f$. Hence $\int_{P_0}f = \int_{P_N}f$. The last statement follows by fixing $z_0 \in U$ and setting $F(z) = \int_\gamma f$ for any path from $z_0$ to $z$, which is independent
of the path since $U$ is simply connected; then $F' = f$ as in Corollary~\ref{complex-analysis:corollary-primitive-convex}.
\end{proof}

\begin{definition}\label{complex-analysis:definition-winding-number}
For a closed path $\gamma$ and $a \notin \gamma^*$, the \emph{winding number} is $n(\gamma, a) = \frac{1}{2\pi i}\int_\gamma\frac{dz}{z - a}$.
\end{definition}

\begin{proposition}\label{complex-analysis:proposition-winding-number}
$n(\gamma, a)$ is an integer; it is constant on connected components of $\CC \setminus \gamma^*$ and zero on the unbounded component. Homotopic closed paths in $\CC \setminus \{a\}$ have the same winding
number about $a$, and for the circle $\gamma(t) = a + re^{2\pi it}$ one has $n(\gamma, a) = 1$.
\end{proposition}

\begin{proof}
Let $\gamma \colon [0, 1] \to \CC$ and $\varphi(t) = \int_0^t\frac{\gamma'(s)}{\gamma(s) - a}\,ds$. Then $(e^{-\varphi}(\gamma - a))' = 0$ on each $C^1$ piece, so $e^{\varphi(1)} = (\gamma(1) - a)/(\gamma(0) - a) = 1$ and
$\varphi(1) \in 2\pi i\ZZ$ by Theorem~\ref{real-analysis:theorem-pi}. The function $a \mapsto n(\gamma, a)$ is continuous on $\CC \setminus \gamma^*$ and integer valued, hence locally constant, and it tends
to $0$ as $|a| \to \infty$. Homotopy invariance is Theorem~\ref{complex-analysis:theorem-cauchy-homotopy} applied to $z \mapsto 1/(z - a)$ on $\CC \setminus \{a\}$, and the circle was computed in
Lemma~\ref{complex-analysis:lemma-path-integral}.
\end{proof}

\section{The Cauchy integral formula}\label{complex-analysis:section-cauchy-formula}

\begin{theorem}\label{complex-analysis:theorem-cauchy-formula}
Let $U$ be convex, $f$ holomorphic on $U$, $\gamma$ a closed path in $U$ and $a \in U \setminus \gamma^*$. Then
\[
n(\gamma, a)f(a) = \frac{1}{2\pi i}\int_\gamma\frac{f(z)}{z - a}\,dz.
\]
In particular, if $\overline{D}(a, r) \subseteq U$ and $\gamma$ is the circle $\partial D(a, r)$ traversed once, then $f(w) = \frac{1}{2\pi i}\int_\gamma\frac{f(z)}{z - w}dz$ for $w \in D(a, r)$.
\end{theorem}

\begin{proof}
Define $g(z) = (f(z) - f(a))/(z - a)$ for $z \neq a$ and $g(a) = f'(a)$. Then $g$ is continuous on $U$ and holomorphic on $U \setminus \{a\}$, so $\int_\gamma g = 0$ by
Corollary~\ref{complex-analysis:corollary-primitive-convex}. For the second statement, note that $n(\gamma, w) = n(\gamma, a) = 1$ for $w \in D(a, r)$ by
Proposition~\ref{complex-analysis:proposition-winding-number}, and apply the formula in the convex set $D(a, r')$ for $r < r' $ with $\overline{D}(a, r') \subseteq U$.
\end{proof}

\begin{theorem}\label{complex-analysis:theorem-analytic}
Let $f$ be holomorphic on an open $U$ and $D(a, R) \subseteq U$. Then $f(w) = \sum_{n \geq 0}c_n(w - a)^n$ on $D(a, R)$, with
$c_n = \frac{1}{2\pi i}\int_{\partial D(a, r)}\frac{f(z)}{(z - a)^{n+1}}dz = \frac{f^{(n)}(a)}{n!}$ for any $0 < r < R$. In particular holomorphic functions have holomorphic derivatives of all orders, and
$|f^{(n)}(a)| \leq n!\,r^{-n}\sup_{|z - a| = r}|f|$ (\emph{Cauchy estimates}).
\end{theorem}

\begin{proof}
Fix $|w - a| < r < R$. For $|z - a| = r$, $\frac{1}{z - w} = \sum_{n \geq 0}\frac{(w - a)^n}{(z - a)^{n+1}}$, converging uniformly in $z$ since $|w - a|/r < 1$. Multiplying by $f(z)/2\pi i$ and integrating term by
term (uniform convergence, Proposition~\ref{real-analysis:proposition-riemann}) gives the expansion with the stated $c_n$, by Theorem~\ref{complex-analysis:theorem-cauchy-formula}. A power series is
infinitely complex differentiable with $c_n = f^{(n)}(a)/n!$ (Theorem~\ref{real-analysis:theorem-power-series-derivative}), and the estimate follows from
Lemma~\ref{complex-analysis:lemma-path-integral}.
\end{proof}

The last statement of Theorem~\ref{complex-analysis:theorem-analytic} deserves a closer look, because nothing like it is true for real functions. Definition~\ref{complex-analysis:definition-holomorphic} only
asks that $f^{\prime}(a)$ exists at every point; it does not ask that $f^{\prime}$ be continuous. For a function of one real variable this would be a much weaker requirement than being of class $C^1$:
$x^2\sin(1/x)$, extended by $0$ at $0$, is differentiable everywhere with a discontinuous derivative, and its antiderivatives are differentiable twice but not three times. For complex differentiability
the hierarchy collapses: one derivative gives all of them, and they are all continuous.

\begin{corollary}\label{complex-analysis:corollary-smooth-harmonic}
Let $f = u + iv$ be holomorphic on an open $U \subseteq \CC$, and identify $\CC$ with $\RR^2$, $z = x + iy$.
\begin{enumerate}
\item The derivatives $f^{\prime}, f^{\prime\prime}, \ldots, f^{(k)}, \ldots$ exist and are holomorphic, hence continuous, on $U$.
\item As a function of the two real variables $x, y$, $f$ is of class $C^\infty$, and for all $a, b \geq 0$
\[
\partial_x^a\partial_y^bf = i^bf^{(a + b)} .
\]
In particular $u$ and $v$ are of class $C^\infty$.
\item $u$ and $v$ are \emph{harmonic}: $\partial_x^2u + \partial_y^2u = 0$ and $\partial_x^2v + \partial_y^2v = 0$ on $U$.
\end{enumerate}
\end{corollary}

\begin{proof}
(1) By Theorem~\ref{complex-analysis:theorem-analytic}, applied on discs $D(a, R) \subseteq U$, the derivative of a holomorphic function is holomorphic. By induction on $k$ every $f^{(k)}$ exists and is
holomorphic. A function that is complex differentiable at a point is continuous there, since $f(a + h) - f(a) = h \cdot (f(a + h) - f(a))/h \to 0 \cdot f^{\prime}(a) = 0$.

(2) Let $g$ be holomorphic on $U$. By Proposition~\ref{complex-analysis:proposition-cauchy-riemann}, $g$ is real differentiable and $Dg(a)$ is multiplication by $g^{\prime}(a)$. The partial derivatives are
the values of $Dg(a)$ on the basis vectors $1$ and $i$ of $\RR^2 = \CC$:
\[
\partial_xg(a) = g^{\prime}(a) \cdot 1 = g^{\prime}(a), \qquad \partial_yg(a) = g^{\prime}(a) \cdot i = ig^{\prime}(a).
\]
So $\partial_x$ and $\partial_y$ map holomorphic functions to holomorphic functions, acting as $g \mapsto g^{\prime}$ and $g \mapsto ig^{\prime}$. Applying this $a + b$ times, starting from $g = f$ and using
(1), shows that every partial derivative $\partial_x^a\partial_y^bf$ (in any order of differentiation) exists and equals $i^bf^{(a + b)}$, which is continuous by (1). Hence $f$ is $C^\infty$ by
definition, and so are its real and imaginary parts $u = \operatorname{Re}f$ and $v = \operatorname{Im}f$, whose partial derivatives are the real and imaginary parts of those of $f$.

(3) By (2), $\partial_x^2f + \partial_y^2f = f^{\prime\prime} + i^2f^{\prime\prime} = 0$. The real and imaginary parts of this equation are $\Delta u = 0$ and $\Delta v = 0$, where
$\Delta = \partial_x^2 + \partial_y^2$, since $\partial_x$ and $\partial_y$ commute with taking real and imaginary parts.
\end{proof}

\begin{remark}\label{complex-analysis:remark-regularity}
\emph{Where does the regularity come from?} The proof of Corollary~\ref{complex-analysis:corollary-smooth-harmonic} passes through four steps, and it is worth seeing what each one uses.
\begin{enumerate}
\item \emph{Goursat's theorem} (Theorem~\ref{complex-analysis:theorem-goursat}) uses complex differentiability only at a single point $z_0$ for each triangle, through the linear approximation
$f(z) = f(z_0) + f^{\prime}(z_0)(z - z_0) + o(|z - z_0|)$. No continuity of $f^{\prime}$ is needed. This is the crucial step: the classical proof of Cauchy's theorem by Green's formula,
$\int_{\partial T}f\,dz = \iint_T(\partial_xf + i\partial_yf)\,i\,dx\,dy$, needs $f$ to be of class $C^1$, and would give the result only for $C^1$ functions.
\item From vanishing integrals over triangles one gets a primitive on discs (Corollary~\ref{complex-analysis:corollary-primitive-convex}), and from that the Cauchy integral formula
(Theorem~\ref{complex-analysis:theorem-cauchy-formula}), applied to the auxiliary function $g(z) = (f(z) - f(a))/(z - a)$, which is continuous at $a$ and holomorphic elsewhere. This is why
Goursat's theorem allows one exceptional point.
\item The Cauchy formula $f(w) = \frac1{2\pi i}\int_{|z - a| = r}\frac{f(z)}{z - w}\,dz$ expresses $f$ inside the disc as an average of its values on the circle against the kernel $1/(z - w)$, which is
infinitely differentiable in $w$. All the regularity of the kernel passes to $f$: one may differentiate under the integral sign, or, as in Theorem~\ref{complex-analysis:theorem-analytic}, expand the kernel
in a geometric series and obtain a power series.
\item Power series are infinitely differentiable (Theorem~\ref{real-analysis:theorem-power-series-derivative}).
\end{enumerate}
So the equation $\partial f/\partial\bar z = 0$ (the Cauchy--Riemann equations, see below) forces smoothness. This is the first instance of \emph{elliptic regularity}: solutions of elliptic
differential equations are as smooth as the equation allows, whatever regularity was assumed a priori; Chapter~\ref{elliptic-operators} proves this in general, and it is what makes harmonic forms smooth in
the Hodge theorem.
\end{remark}

It is often convenient to package the two real partial derivatives into the \emph{Wirtinger derivatives}
\[
\frac{\partial}{\partial z} = \frac12\Bigl(\partial_x - i\partial_y\Bigr), \qquad \frac{\partial}{\partial\bar z} = \frac12\Bigl(\partial_x + i\partial_y\Bigr),
\]
defined for real differentiable functions. For $f = u + iv$ the real and imaginary parts of $2\,\partial f/\partial\bar z = (\partial_xu - \partial_yv) + i(\partial_xv + \partial_yu)$ are the two
Cauchy--Riemann expressions, so a real differentiable $f$ is holomorphic exactly when $\partial f/\partial\bar z = 0$, and then $\partial f/\partial z = f^{\prime}$. On $C^2$ functions
\[
4\frac{\partial}{\partial z}\frac{\partial}{\partial\bar z} = (\partial_x - i\partial_y)(\partial_x + i\partial_y) = \partial_x^2 + \partial_y^2 = \Delta,
\]
since $\partial_x$ and $\partial_y$ commute on $C^2$ functions (Theorem~\ref{real-analysis:theorem-schwarz}). This gives a second proof of part (3): $\Delta f = 4\frac{\partial}{\partial z}\bigl(\frac{\partial f}{\partial\bar z}\bigr) = 0$.
The Wirtinger derivatives satisfy the product rule, and $\partial\bar g/\partial\bar z = \overline{\partial g/\partial z}$, as one checks from the definitions.

\begin{example}\label{complex-analysis:example-harmonic}
\begin{enumerate}
\item From $z^2 = (x^2 - y^2) + i\,2xy$, the functions $x^2 - y^2$ and $2xy$ are harmonic; from $e^z = e^x\cos y + ie^x\sin y$, so are $e^x\cos y$ and $e^x\sin y$. The level curves of $u$ and
$v$ meet at right angles wherever $f^{\prime} \neq 0$: by the Cauchy--Riemann equations, $\nabla u = (\partial_xu, \partial_yu)$ and $\nabla v = (-\partial_yu, \partial_xu)$ are orthogonal vectors of
the same length $|f^{\prime}|$.
\item \emph{$\log|f|$ is harmonic where $f \neq 0$.} On a disc $D \subseteq U$ on which $f$ has no zeros, $f^{\prime}/f$ is holomorphic, so it has a primitive $g$
(Corollary~\ref{complex-analysis:corollary-primitive-convex}); then $(fe^{-g})^{\prime} = (f^{\prime} - fg^{\prime})e^{-g} = 0$, so $f = ce^g$ with a constant $c = e^{c_0} \neq 0$, and
$\log|f| = \operatorname{Re}(g + c_0)$ is harmonic. For example $\log|z| = \frac12\log(x^2 + y^2)$ is harmonic on $\CC^\times$, although it is not the real part of a holomorphic function $g$ on all of
$\CC^\times$: by the Cauchy--Riemann equations $g^{\prime} = \partial_x\operatorname{Re}g - i\partial_y\operatorname{Re}g$, which for $\operatorname{Re}g = \log|z|$ is $(x - iy)/(x^2 + y^2) = 1/z$, and
$1/z$ has no primitive on $\CC^\times$, its integral over the unit circle being $2\pi i$ (Lemma~\ref{complex-analysis:lemma-path-integral}).
\item \emph{$|f|^2$ is not harmonic, but subharmonic.} With the product rule, $\frac{\partial}{\partial\bar z}(f\bar f) = f\,\overline{f^{\prime}}$, because $\partial f/\partial\bar z = 0$ and
$\partial\bar f/\partial\bar z = \overline{f^{\prime}}$; differentiating again, $\frac{\partial}{\partial z}\bigl(f\,\overline{f^{\prime}}\bigr) = f^{\prime}\,\overline{f^{\prime}}$, since $\overline{f^{\prime}}$
has $\partial/\partial z$ equal to $\overline{\partial f^{\prime}/\partial\bar z} = 0$. Hence $\Delta|f|^2 = 4|f^{\prime}|^2 \geq 0$. For $f(z) = z$ this is $\Delta(x^2 + y^2) = 4$.
\end{enumerate}
\end{example}

\begin{counterexample}\label{complex-analysis:counterexample-cauchy-riemann-point}
The Cauchy--Riemann equations at a point do not imply complex differentiability there; Proposition~\ref{complex-analysis:proposition-cauchy-riemann} also requires real differentiability. Let
$f(x + iy) = \sqrt{|xy|}$, so $u = \sqrt{|xy|}$ and $v = 0$. Since $u$ vanishes on both axes, $\partial_xu(0) = \partial_yu(0) = 0$, and the Cauchy--Riemann equations hold at $0$. But for $h = t(1 + i)$
with $t \in \RR^\times$, $(f(h) - f(0))/h = |t|/(t(1 + i))$, which is $\frac1{1 + i}$ for $t > 0$ and $-\frac1{1 + i}$ for $t < 0$; so $f^{\prime}(0)$ does not exist. Existence of the partial derivatives
only concerns the two directions of the axes, while complex differentiability concerns all directions at once.
\end{counterexample}

\begin{corollary}\label{complex-analysis:corollary-liouville}
\begin{enumerate}
\item (Liouville) A bounded entire function is constant; more generally, if $|f(z)| \leq C(1 + |z|)^m$ for an entire $f$, then $f$ is a polynomial of degree at most $m$.
\item (Fundamental theorem of algebra) Every nonconstant polynomial over $\CC$ has a root (cf. Chapter~\ref{fields}).
\item (Morera) A continuous function on an open set whose integral over the boundary of every closed triangle contained in the set vanishes is holomorphic.
\item (Weierstrass) If holomorphic functions $f_n$ on $U$ converge uniformly on compact subsets to $f$, then $f$ is holomorphic and $f_n^{(k)} \to f^{(k)}$ uniformly on compact subsets.
\end{enumerate}
\end{corollary}

\begin{proof}
(1) By the Cauchy estimates with $r \to \infty$, $f^{(n)}(a) = 0$ for $n > m$ and all $a$. (2) If $p$ has no root, $1/p$ is entire and tends to $0$ at infinity, hence is bounded and constant. (3) Locally on
a disc the hypothesis gives a primitive $F$ as in Corollary~\ref{complex-analysis:corollary-primitive-convex}, and $F$ is holomorphic, hence so is $F' = f$ by
Theorem~\ref{complex-analysis:theorem-analytic}. (4) The uniform limit is continuous and satisfies Morera's condition by Theorem~\ref{complex-analysis:theorem-goursat} and uniform convergence; convergence
of derivatives follows from the Cauchy integral formula for $f^{(k)}$ on circles inside $U$.
\end{proof}

\begin{remark}\label{complex-analysis:remark-dbar-hypotheses}
\emph{Is $f$ holomorphic if and only if $\partial f/\partial\bar z = 0$?} Yes, with two provisos: the equation must hold on an open set, and $f$ must be real differentiable there (for instance of class $C^1$), since
the Wirtinger derivatives are only defined for such functions.
\begin{enumerate}
\item \emph{Open sets.} Holomorphy is a condition on an open set, and $\partial f/\partial\bar z = 0$ at a single point only gives complex differentiability at that point. For $f(z) = |z|^2 = z\bar z$,
$\partial f/\partial\bar z = z$ vanishes only at $0$: $f$ is complex differentiable at $0$, with $f^{\prime}(0) = \lim\bar h = 0$, and nowhere else, so it is holomorphic on no open set.
\item \emph{Real differentiability.} Requiring only that the partial derivatives exist and satisfy the Cauchy--Riemann equations is not enough, even at every point of an open set. Let
$f(z) = e^{-z^{-4}}$ for $z \neq 0$ and $f(0) = 0$. Away from $0$, $f$ is holomorphic. On the real axis $f(x) = e^{-x^{-4}}$, and on the imaginary axis $f(iy) = e^{-y^{-4}}$ since $(iy)^4 = y^4$; both
tend to $0$ faster than any power, so $\partial_xf(0) = \partial_yf(0) = 0$ and the Cauchy--Riemann equations hold everywhere. But on the diagonal $z = re^{i\pi/4}$, $z^4 = -r^4$ and
$f(z) = e^{r^{-4}} \to \infty$ as $r \to 0$: $f$ is not even continuous at $0$. Counterexample~\ref{complex-analysis:counterexample-cauchy-riemann-point} shows the same phenomenon at a single point
for a continuous function.
\item \emph{The Looman--Menchoff theorem.} Continuity is the only thing missing in (2): if $f$ is continuous on $U$, its partial derivatives exist at every point of $U$ and satisfy the Cauchy--Riemann equations,
then $f$ is holomorphic \cite{gray-morris-1978}. The proof is delicate (it uses Baire's theorem); we do not need it.
\item \emph{Weak solutions.} The equation can also be imposed in the weak sense, by moving the derivative onto test functions; then continuity alone suffices, as the next proposition shows. This is the point of view of
distributions (Chapter~\ref{distributions}), in which every locally integrable weak solution of $\partial f/\partial\bar z = 0$ agrees almost everywhere with a holomorphic function, a first case of elliptic
regularity (Remark~\ref{complex-analysis:remark-regularity}).
\end{enumerate}
\end{remark}

\begin{proposition}\label{complex-analysis:proposition-weak-cauchy-riemann}
Let $U \subseteq \CC$ be open and $f \colon U \to \CC$ continuous. Then $f$ is holomorphic if and only if
\[
\int_Uf\,\frac{\partial\varphi}{\partial\bar z}\,dA = 0 \qquad \text{for every smooth } \varphi \colon U \to \CC \text{ with compact support,}
\]
where $dA = dx\,dy$ is Lebesgue measure.
\end{proposition}

\begin{proof}
\emph{Necessity.} If $f$ is holomorphic, it is $C^\infty$ (Corollary~\ref{complex-analysis:corollary-smooth-harmonic}). Extend $f\varphi$ by $0$ outside the support of $\varphi$; it is then a $C^1$ function on
$\CC$ with compact support. For fixed $y$, $\int_\RR\partial_x(f\varphi)\,dx = 0$ by the fundamental theorem of calculus, and likewise for $\partial_y$; integrating over the other variable (Fubini's theorem),
$\int(f\,\partial_x\varphi + \varphi\,\partial_xf)\,dA = 0$, and similarly for $y$. Combining, $\int f\,\partial\varphi/\partial\bar z\,dA = -\int\varphi\,\partial f/\partial\bar z\,dA = 0$.

\emph{Sufficiency.} By Corollary~\ref{real-analysis:corollary-bump-function} there is a smooth $\rho \geq 0$ on $\CC$ with support in the unit disc and $\int\rho\,dA = 1$; put $\rho_\varepsilon(w) = \varepsilon^{-2}\rho(w/\varepsilon)$,
which has support in $\overline{D}(0, \varepsilon)$ and integral $1$. Let $U_\varepsilon = \{z : \overline{D}(z, \varepsilon) \subseteq U\}$, an open set, and for $z \in U_\varepsilon$
\[
f_\varepsilon(z) = \int_Uf(w)\,\rho_\varepsilon(z - w)\,dA(w).
\]
\emph{Step 1: $f_\varepsilon$ is $C^1$ on $U_\varepsilon$, with derivatives obtained by differentiating $\rho_\varepsilon$.} Fix $z_0 \in U_\varepsilon$ and $\delta > 0$ with $\overline{D}(z_0, \varepsilon + \delta) \subseteq U$.
For $|z - z_0| < \delta$ the integrand vanishes outside the compact disc $K = \overline{D}(z_0, \varepsilon + \delta)$, and $|f(w)\,\partial_x\rho_\varepsilon(z - w)| \leq \sup_K|f| \cdot \sup|\partial_x\rho_\varepsilon| \cdot 1_K(w)$,
an integrable bound independent of $z$. By Proposition~\ref{measure-theory:proposition-differentiation-integral}, applied to the real and imaginary parts, $\partial_xf_\varepsilon(z) = \int f(w)(\partial_x\rho_\varepsilon)(z - w)\,dA(w)$,
and similarly for $\partial_y$. These are continuous in $z$, by dominated convergence with the same bound, so $f_\varepsilon$ is $C^1$, and
\[
\frac{\partial f_\varepsilon}{\partial\bar z}(z) = \int_Uf(w)\,\frac{\partial\rho_\varepsilon}{\partial\bar z}(z - w)\,dA(w).
\]
\emph{Step 2: $f_\varepsilon$ is holomorphic.} For fixed $z \in U_\varepsilon$ the function $\varphi(w) = \rho_\varepsilon(z - w)$ is smooth with compact support in $U$, and by the chain rule
$\partial_x\varphi(w) = -(\partial_x\rho_\varepsilon)(z - w)$ and $\partial_y\varphi(w) = -(\partial_y\rho_\varepsilon)(z - w)$, so $\partial\varphi/\partial\bar w = -(\partial\rho_\varepsilon/\partial\bar z)(z - w)$. The hypothesis gives
$\partial f_\varepsilon/\partial\bar z(z) = -\int f\,\partial\varphi/\partial\bar w\,dA = 0$. So $f_\varepsilon$ is real differentiable with $\partial f_\varepsilon/\partial\bar z = 0$ on $U_\varepsilon$, hence holomorphic there.

\emph{Step 3: $f_\varepsilon \to f$ uniformly on compact sets.} Let $C \subseteq U$ be compact and $\varepsilon_0 > 0$ with $C^{\prime} = \{w : \operatorname{dist}(w, C) \leq \varepsilon_0\} \subseteq U$; then $C \subseteq U_\varepsilon$
for $\varepsilon \leq \varepsilon_0$. Substituting $w \mapsto z - w$ and using $\int\rho_\varepsilon = 1$, for $z \in C$
\[
|f_\varepsilon(z) - f(z)| = \Bigl|\int\bigl(f(z - w) - f(z)\bigr)\rho_\varepsilon(w)\,dA(w)\Bigr| \leq \sup\{|f(z - w) - f(z)| : |w| \leq \varepsilon\},
\]
which tends to $0$ uniformly in $z \in C$, because $f$ is uniformly continuous on the compact set $C^{\prime}$. Every compact subset of $U$ lies in $U_\varepsilon$ for small $\varepsilon$, so by Weierstrass's theorem
(Corollary~\ref{complex-analysis:corollary-liouville}(4)), applied on the open sets $U_{\varepsilon_0}$, $f$ is holomorphic on each $U_{\varepsilon_0}$, hence on $U$.
\end{proof}

The function $f_\varepsilon$ is a \emph{regularisation} of $f$ by convolution with the \emph{mollifier} $\rho_\varepsilon$: it is smooth, close to $f$, and it inherits every linear constant-coefficient equation that $f$
satisfies weakly. The weak equation must hold on all of $U$: $1/z$ is holomorphic on $\CC^\times$ and locally integrable on $\CC$, but on $\CC$ it does not satisfy the weak equation, since
$\int(1/z)\,\partial\varphi/\partial\bar z\,dA = -\pi\varphi(0)$ (Exercise~\ref{complex-analysis:exercise-cauchy-kernel-dbar}); the point $0$ carries a ``source''.

\begin{theorem}[Identity theorem]\label{complex-analysis:theorem-identity}
Let $f$ be holomorphic on a domain $U$. If the zero set of $f$ has a limit point in $U$, then $f \equiv 0$. Consequently two holomorphic functions agreeing on a set with a limit point in $U$ are equal, and
the zeros of a nonzero $f$ are isolated, each of finite order $m$ with $f(z) = (z - a)^mg(z)$, $g$ holomorphic and $g(a) \neq 0$.
\end{theorem}

\begin{proof}
Let $A$ be the set of $a \in U$ with $f^{(n)}(a) = 0$ for all $n$; it is closed in $U$, and open because $f$ vanishes identically near such a point by
Theorem~\ref{complex-analysis:theorem-analytic}. If $z_k \to a$ with $z_k \neq a$ and $f(z_k) = 0$, and if some derivative at $a$ were nonzero, then with $m$ minimal, $f(z) = (z - a)^mg(z)$ with $g(a) \neq 0$
by the power series expansion, so $f$ would have no zeros near $a$ other than $a$. Hence $a \in A$, and $A = U$ by connectedness.
\end{proof}

\begin{theorem}[Maximum modulus principle]\label{complex-analysis:theorem-maximum-modulus}
Let $f$ be holomorphic on a domain $U$. If $|f|$ has a local maximum in $U$, then $f$ is constant. If moreover $U$ is bounded and $f$ extends continuously to $\overline{U}$, then $|f|$ attains its maximum on
$\partial U$.
\end{theorem}

\begin{proof}
Theorem~\ref{complex-analysis:theorem-cauchy-formula} on circles gives the mean value property $f(a) = \int_0^1f(a + re^{2\pi it})\,dt$. If $|f| \leq |f(a)|$ on $D(a, R)$, then for $r < R$,
$|f(a)| \leq \int_0^1|f(a + re^{2\pi it})|\,dt \leq |f(a)|$, so $|f| = |f(a)|$ on every such circle by continuity, and $|f|$ is constant on $D(a, R)$. Then $f\bar f$ is constant; if $f(a) \neq 0$ this makes
$\bar f = c/f$ holomorphic, so both $f$ and $\bar f$ satisfy the Cauchy--Riemann equations and $f' = 0$ on the disc. In either case $f$ is constant on $D(a, R)$, hence on $U$ by
Theorem~\ref{complex-analysis:theorem-identity}. The last statement follows from compactness of $\overline{U}$.
\end{proof}

\section{Singularities, Laurent series and residues}\label{complex-analysis:section-residues}

\begin{theorem}[Laurent expansion]\label{complex-analysis:theorem-laurent}
Let $f$ be holomorphic on the annulus $A = \{r < |z - a| < R\}$, $0 \leq r < R \leq \infty$. There are unique $c_n \in \CC$, $n \in \ZZ$, with $f(z) = \sum_{n \in \ZZ}c_n(z - a)^n$, converging absolutely and
uniformly on compact subsets of $A$, and $c_n = \frac{1}{2\pi i}\int_{\partial D(a, \rho)}\frac{f(z)}{(z - a)^{n+1}}dz$ for any $r < \rho < R$.
\end{theorem}

\begin{proof}
Fix $w$ with $r < r_1 < |w - a| < r_2 < R$. The function $g(z) = (f(z) - f(w))/(z - w)$, extended by $f'(w)$ at $w$, is continuous on $A$ and holomorphic off $w$, hence holomorphic by
Corollary~\ref{complex-analysis:corollary-liouville}(3) and Theorem~\ref{complex-analysis:theorem-goursat}. The circles $\partial D(a, r_1)$ and $\partial D(a, r_2)$ are homotopic as closed paths in $A$ (move
radially), so $\int_{\partial D(a, r_2)}g = \int_{\partial D(a, r_1)}g$ by Theorem~\ref{complex-analysis:theorem-cauchy-homotopy}. Since $n(\partial D(a, r_2), w) = 1$ and $n(\partial D(a, r_1), w) = 0$, this
gives
\[
f(w) = \frac{1}{2\pi i}\int_{\partial D(a, r_2)}\frac{f(z)}{z - w}dz - \frac{1}{2\pi i}\int_{\partial D(a, r_1)}\frac{f(z)}{z - w}dz.
\]
In the first integral expand $\frac{1}{z - w} = \sum_{n \geq 0}\frac{(w - a)^n}{(z - a)^{n+1}}$, in the second $\frac{-1}{z - w} = \sum_{n \geq 1}\frac{(z - a)^{n-1}}{(w - a)^n}$, both uniformly convergent on the
respective circles, and integrate term by term. The resulting coefficients are independent of the radius, again by homotopy invariance. Uniqueness follows by integrating the series against
$(z - a)^{-n-1}$ term by term.
\end{proof}

\begin{definition}\label{complex-analysis:definition-singularity}
Let $f$ be holomorphic on a punctured disc $D(a, R) \setminus \{a\}$, with Laurent coefficients $c_n$. The singularity at $a$ is \emph{removable} if $c_n = 0$ for all $n < 0$, a \emph{pole of order $m$} if
$c_{-m} \neq 0$ and $c_n = 0$ for $n < -m$, and \emph{essential} otherwise. The \emph{residue} is $\operatorname{Res}_af = c_{-1}$. A function holomorphic on $U$ except for poles at a discrete closed subset
is \emph{meromorphic}; the \emph{order} $\operatorname{ord}_af$ is $m$ at a zero of order $m$ and $-m$ at a pole of order $m$.
\end{definition}

\begin{proposition}\label{complex-analysis:proposition-singularities}
Let $f$ be holomorphic on $D(a, R) \setminus \{a\}$.
\begin{enumerate}
\item (Riemann) If $f$ is bounded near $a$, the singularity is removable.
\item The singularity is a pole if and only if $|f(z)| \to \infty$ as $z \to a$.
\item (Casorati--Weierstrass) If the singularity is essential, $f(D(a, \varepsilon) \setminus \{a\})$ is dense in $\CC$ for every $\varepsilon > 0$.
\end{enumerate}
\end{proposition}

\begin{proof}
(1) If $|f| \leq M$ near $a$, then $|c_{-n}| \leq M\rho^{n}$ for small $\rho$ by Theorem~\ref{complex-analysis:theorem-laurent}, so $c_{-n} = 0$ for $n \geq 1$. (2) If $f$ has a pole of order $m$, then
$(z - a)^mf(z)$ extends holomorphically with nonzero value at $a$, so $|f| \to \infty$. Conversely, if $|f| \to \infty$, then $1/f$ is holomorphic and nonvanishing near $a$ with limit $0$, so by (1) it extends
with a zero of some order $m$ at $a$, and $f = 1/(1/f)$ has a pole of order $m$. (3) If the image of some punctured disc misses $D(w, \delta)$, then $g = 1/(f - w)$ is bounded there, hence extends
holomorphically by (1); then $f = w + 1/g$ has a removable singularity or a pole by (2).
\end{proof}

\begin{example}\label{complex-analysis:example-essential}
The function $f(z) = e^{1/z}$ on $\CC \setminus \{0\}$ has Laurent series $\sum_{n \geq 0}z^{-n}/n!$, so $0$ is an essential singularity. More than Casorati--Weierstrass holds here: every $w \neq 0$ is taken
infinitely often in every punctured disc around $0$. Indeed, if $w = e^c$, the points $z_k = 1/(c + 2\pi ik)$, $k \in \ZZ$, satisfy $f(z_k) = w$ and tend to $0$. The value $0$ is never taken. On the real axis,
$e^{1/x} \to \infty$ as $x \to 0^+$ and $\to 0$ as $x \to 0^-$, while along the imaginary axis $|e^{1/(iy)}| = 1$: the behaviour near an essential singularity depends on the direction of approach.
\end{example}

\begin{counterexample}\label{complex-analysis:counterexample-real-radius}
Real analysis does not see the radius of convergence. The function $g(x) = 1/(1 + x^2)$ is real-analytic on all of $\RR$, but its Taylor series $\sum_n(-1)^nx^{2n}$ at $0$ has radius of convergence $1$. The reason
is complex: $g$ extends to the meromorphic function $1/(1 + z^2)$ with poles at $\pm i$, at distance $1$ from $0$. By Theorem~\ref{complex-analysis:theorem-analytic} the series converges on $D(0, 1)$, and it cannot converge on a
larger disc: its sum would be holomorphic there, equal to $1/(1 + z^2)$ on $D(0, 1)$ by the identity theorem, hence bounded near $i$, whereas $|1/(1 + z^2)| \to \infty$ as $z \to i$. Similarly $x \mapsto e^{-1/x^2}$ (extended by $0$) is smooth on $\RR$ with all derivatives zero at $0$, so it is not real-analytic there; as a complex function,
$e^{-1/z^2}$ has an essential singularity at $0$, and it is unbounded along the imaginary axis.
\end{counterexample}

\begin{lemma}\label{complex-analysis:lemma-difference-quotient}
Let $h$ be holomorphic on an open $U$ and define $g \colon U \times U \to \CC$ by $g(z, w) = (h(z) - h(w))/(z - w)$ for $z \neq w$ and $g(z, z) = h'(z)$. Then $g$ is continuous, and holomorphic in each
variable separately.
\end{lemma}

\begin{proof}
Holomorphy in each variable off the diagonal is clear, and at $z = w$ the singularity is removable by Proposition~\ref{complex-analysis:proposition-singularities}(1), since $g$ is bounded near it by the
estimate below. For continuity at $(a, a)$, take $\overline{D}(a, r) \subseteq U$; for $z, w \in D(a, r)$ the segment $[w, z]$ lies in the disc and
$g(z, w) = \int_0^1h'(w + t(z - w))\,dt$ (by the fundamental theorem of calculus applied to $t \mapsto h(w + t(z - w))$, also valid when $z = w$). Since $h'$ is continuous, hence uniformly continuous on
$\overline{D}(a, r/2)$, this expression is continuous in $(z, w)$ there.
\end{proof}

\begin{theorem}[Cauchy's theorem, homology form]\label{complex-analysis:theorem-cauchy-homology}
Let $h$ be holomorphic on an open $U \subseteq \CC$ and let $\gamma$ be a closed path in $U$ with $n(\gamma, z) = 0$ for every $z \in \CC \setminus U$. Then $\int_\gamma h = 0$, and
$n(\gamma, w)h(w) = \frac{1}{2\pi i}\int_\gamma\frac{h(z)}{z - w}\,dz$ for every $w \in U \setminus \gamma^*$.
\end{theorem}

\begin{proof}
(Dixon.) With $g$ as in Lemma~\ref{complex-analysis:lemma-difference-quotient}, let $\varphi(w) = \frac{1}{2\pi i}\int_\gamma g(z, w)\,dz$ for $w \in U$. It is continuous, since $g$ is uniformly continuous on
$\gamma^* \times K$ for compact $K \subseteq U$, and holomorphic: for a closed triangle $T \subseteq U$, Fubini's theorem for continuous functions
(Proposition~\ref{real-analysis:proposition-parameter-integrals}) and Theorem~\ref{complex-analysis:theorem-goursat} give
$\int_{\partial T}\varphi(w)\,dw = \frac{1}{2\pi i}\int_\gamma\left(\int_{\partial T}g(z, w)\,dw\right)dz = 0$, so Morera's theorem applies. Let $V = \{w \in \CC \setminus \gamma^* : n(\gamma, w) = 0\}$, an open set,
and $\psi(w) = \frac{1}{2\pi i}\int_\gamma\frac{h(z)}{z - w}\,dz$ on $V$, which is holomorphic (same argument). For $w \in U \cap V$,
$\varphi(w) = \psi(w) - n(\gamma, w)h(w) = \psi(w)$. By hypothesis $U \cup V = \CC$, so $\varphi$ and $\psi$ define an entire function $F$. The unbounded component of $\CC \setminus \gamma^*$ lies in $V$, and
$\psi(w) \to 0$ as $|w| \to \infty$; so $F$ is bounded, hence constant equal to $0$ by Liouville's theorem (Corollary~\ref{complex-analysis:corollary-liouville}). Thus $\varphi \equiv 0$, which is the integral
formula. Applying it to $z \mapsto (z - w)h(z)$ at the point $w$ gives $0 = \frac{1}{2\pi i}\int_\gamma h$.
\end{proof}

\begin{theorem}[Residue theorem]\label{complex-analysis:theorem-residues}
Let $U$ be simply connected and open, $S \subseteq U$ discrete and closed in $U$, and $f$ holomorphic on $U \setminus S$. For every closed path $\gamma$ in $U \setminus S$,
\[
\frac{1}{2\pi i}\int_\gamma f = \sum_{a \in S}n(\gamma, a)\operatorname{Res}_af,
\]
a finite sum.
\end{theorem}

\begin{proof}
Let $W = \{z \in \CC \setminus \gamma^* : n(\gamma, z) \neq 0\}$, a union of bounded components of $\CC \setminus \gamma^*$, so $K = \overline{W} \subseteq W \cup \gamma^*$ is compact. If $z \notin U$, then $1/(w - z)$ is
holomorphic on the simply connected $U$, hence has a primitive, so $n(\gamma, z) = 0$; therefore $K \subseteq U$, and $S \cap K$ is finite, say $\{a_1, \ldots, a_k\}$, because $S$ is discrete and closed in $U$.
Let $P_j(z) = \sum_{n \geq 1}c^{(j)}_{-n}(z - a_j)^{-n}$ be the principal part of $f$ at $a_j$. It converges locally uniformly on $\CC \setminus \{a_j\}$, because the Laurent series of $f$ converges on the
punctured disc and the principal part, as a power series in $1/(z - a_j)$, then has infinite radius of convergence; so $P_j$ is holomorphic there. The function $h = f - \sum_jP_j$ has removable
singularities at $a_1, \ldots, a_k$, hence extends holomorphically to $U' = U \setminus (S \setminus \{a_1, \ldots, a_k\})$. For $z \notin U'$ we have $n(\gamma, z) = 0$: either $z \notin U$, or
$z \in S \setminus \{a_1, \ldots, a_k\}$, which was excluded from $K$. So $\int_\gamma h = 0$ by Theorem~\ref{complex-analysis:theorem-cauchy-homology}. Finally
$\int_\gamma P_j = 2\pi i\,n(\gamma, a_j)c^{(j)}_{-1}$, since the terms with $n \geq 2$ have primitives, and $n(\gamma, a) = 0$ for the remaining $a \in S$.
\end{proof}

\begin{example}\label{complex-analysis:example-residue-computation}
$\int_{-\infty}^{\infty}\frac{dx}{1 + x^2} = \pi$. Integrate $f(z) = 1/(1 + z^2)$ over the boundary of the half-disc $\{|z| \leq R, \operatorname{Im}z \geq 0\}$: the only singularity inside is the simple pole
at $i$ with residue $\lim_{z \to i}(z - i)f(z) = 1/(2i)$, so the integral is $\pi$. The contribution of the semicircle is at most $\pi R/(R^2 - 1) \to 0$, and the integral over $[-R, R]$ tends to the
improper integral.
\end{example}

\begin{theorem}[Argument principle and Rouch\'e]\label{complex-analysis:theorem-argument-principle}
Let $U$ be simply connected, $f$ meromorphic on $U$ and $\gamma$ a closed path in $U$ avoiding the zeros and poles of $f$. Then
$\frac{1}{2\pi i}\int_\gamma\frac{f'}{f} = \sum_a n(\gamma, a)\operatorname{ord}_af$. If $g$ is holomorphic on $U$ and $|f - g| < |f|$ on $\gamma^*$, with $f$ holomorphic, then $f$ and $g$ have the same number of
zeros counted with multiplicity and weighted by $n(\gamma, -)$.
\end{theorem}

\begin{proof}
Near a zero or pole, $f(z) = (z - a)^{m}h(z)$ with $h(a) \neq 0$ and $m = \operatorname{ord}_af$, so $f'/f = m/(z - a) + h'/h$ has residue $m$ at $a$; apply
Theorem~\ref{complex-analysis:theorem-residues}. For Rouch\'e's theorem, the functions $f_t = f + t(g - f)$, $t \in [0, 1]$, have no zeros on $\gamma^*$ because $|f_t| \geq |f| - |g - f| > 0$; the integer
$N(t) = \frac{1}{2\pi i}\int_\gamma f_t'/f_t$ depends continuously on $t$ (the integrand is continuous in $(t, z)$ and $\gamma^*$ is compact), hence is constant.
\end{proof}

\begin{example}\label{complex-analysis:example-fta-proofs}
\emph{More proofs of the fundamental theorem of algebra.} Let $p(z) = z^n + a_{n-1}z^{n-1} + \cdots + a_0$, $n \geq 1$. Besides the algebraic proof (Theorem~\ref{fields:theorem-fta}, by Galois theory and the
intermediate value theorem) and Liouville's theorem (Corollary~\ref{complex-analysis:corollary-liouville}(2)):
\begin{enumerate}
\item \emph{By Rouch\'e's theorem, or by winding numbers.} On the circle $|z| = R$ with $R > 1 + \sum_i|a_i|$, $|p(z) - z^n| \leq \sum_i|a_i|R^i < R^n = |z^n|$. By the theorem, $p$ has as many zeros in
$|z| < R$ as $z^n$, namely $n$, counted with multiplicity. Topologically: the loops $t \mapsto p(Re^{it})$ and $t \mapsto R^ne^{int}$ in $\CC \setminus 0$ are homotopic, so they have the same winding
number $n$ about $0$; if $p$ had no zeros, $t \mapsto p(re^{it})$, $0 \leq r \leq R$, would be a homotopy in $\CC \setminus 0$ to a constant loop, of winding number $0$.
\item \emph{By the minimum modulus (d'Alembert--Argand).} Since $|p(z)| \geq |z|^n(1 - \sum_i|a_i||z|^{i-n}) \to \infty$, $|p|$ attains a minimum on $\CC$ (on a large closed disc by compactness, and outside it
$|p| > |p(0)|$). Suppose the minimum is attained at $z_0$ with $p(z_0) \neq 0$. Expanding in $h$, $p(z_0 + h) = p(z_0) + ch^k + h^{k+1}q(h)$ with $c \neq 0$, $k \geq 1$ and $q$ a polynomial. Choose
$\zeta$ with $c\zeta^k = -p(z_0)$ ($k$-th roots exist by the polar form). For $0 < t < 1$, $p(z_0 + t\zeta) = (1 - t^k)p(z_0) + t^{k+1}\zeta^{k+1}q(t\zeta)$, so
$|p(z_0 + t\zeta)| \leq (1 - t^k)|p(z_0)| + Ct^{k+1} < |p(z_0)|$ for small $t > 0$, contradicting minimality. This proof uses only that $\CC$ has $k$-th roots and that continuous functions on compact sets have
minima.
\end{enumerate}
\end{example}

\begin{theorem}[Open mapping theorem]\label{complex-analysis:theorem-open-mapping}
A nonconstant holomorphic function on a domain is an open map. If it is injective, its derivative has no zeros and its inverse is holomorphic.
\end{theorem}

\begin{proof}
Let $w_0 = f(z_0)$. By Theorem~\ref{complex-analysis:theorem-identity} the zero of $f - w_0$ at $z_0$ has finite order $m \geq 1$ and is isolated, so there is $r$ with $|f - w_0| \geq \delta > 0$ on
$|z - z_0| = r$ and no other zero in $\overline{D}(z_0, r)$. For $|w - w_0| < \delta$, Rouch\'e's theorem applied to $f - w_0$ and $f - w$ on the circle shows that $f - w$ has $m$ zeros in $D(z_0, r)$; so
$D(w_0, \delta) \subseteq f(D(z_0, r))$. If $f$ is injective, then $m = 1$ at every point, that is, $f' \neq 0$; and $f^{-1}$ is continuous because $f$ is open, and differentiable with derivative $1/f'(f^{-1})$ by
the computation of Proposition~\ref{real-analysis:proposition-inverse-one}, valid over $\CC$.
\end{proof}

\section{Conformal maps and the Riemann mapping theorem}\label{complex-analysis:section-riemann-mapping}

\begin{lemma}[Schwarz]\label{complex-analysis:lemma-schwarz}
Let $f \colon \mathbb{D} \to \mathbb{D}$ be holomorphic with $f(0) = 0$. Then $|f(z)| \leq |z|$ and $|f'(0)| \leq 1$, with equality for some $z \neq 0$ or in the second inequality only if $f(z) = \lambda z$ with
$|\lambda| = 1$.
\end{lemma}

\begin{proof}
The function $g(z) = f(z)/z$, extended by $f'(0)$ at $0$, is holomorphic on $\mathbb{D}$ (Proposition~\ref{complex-analysis:proposition-singularities}(1)). On $|z| = r$, $|g| \leq 1/r$, so by the maximum
principle $|g| \leq 1/r$ on $D(0, r)$; letting $r \to 1$ gives $|g| \leq 1$. Equality at an interior point forces $g$ to be a constant of modulus $1$.
\end{proof}

\begin{proposition}\label{complex-analysis:proposition-disc-automorphisms}
For $a \in \mathbb{D}$ the M\"obius map $\phi_a(z) = \frac{z - a}{1 - \bar az}$ is a biholomorphic self-map of $\mathbb{D}$ with $\phi_a(a) = 0$ and inverse $\phi_{-a}$. Every biholomorphic self-map of
$\mathbb{D}$ is $z \mapsto \lambda\phi_a(z)$ with $|\lambda| = 1$.
\end{proposition}

\begin{proof}
For $|z| < 1$ and $|a| < 1$, $|1 - \bar az|^2 - |z - a|^2 = (1 - |a|^2)(1 - |z|^2) > 0$, so $\phi_a$ maps $\mathbb{D}$ into $\mathbb{D}$, and a computation gives $\phi_{-a}\phi_a = \id$. If $F$ is a biholomorphic
self-map with $F(a) = 0$, then $g = F \circ \phi_{-a}$ fixes $0$, and Lemma~\ref{complex-analysis:lemma-schwarz} applied to $g$ and $g^{-1}$ gives $|g'(0)| = 1$, hence $g(z) = \lambda z$.
\end{proof}

\begin{theorem}[Montel]\label{complex-analysis:theorem-montel}
Let $\mathcal{F}$ be a family of holomorphic functions on an open $U$ that is uniformly bounded on every compact subset. Then every sequence in $\mathcal{F}$ has a subsequence converging uniformly on compact
subsets to a holomorphic function.
\end{theorem}

\begin{proof}
Let $K \subseteq U$ be compact and $2\rho$ the distance from $K$ to $\CC \setminus U$ (or $1$ if $U = \CC$), and let $K'$ be the set of points at distance $\leq \rho$ from $K$, a compact subset of $U$. By the
Cauchy estimates, $|f'| \leq \rho^{-1}\sup_{K'}|f|$ on $K$ for all $f \in \mathcal{F}$, so $\mathcal{F}$ is equicontinuous on $K$ (mean value inequality). By the Arzel\`a--Ascoli theorem
(Theorem~\ref{real-analysis:theorem-arzela-ascoli}) a sequence has a subsequence converging uniformly on $K$. Exhausting $U$ by compact sets $K_1 \subseteq K_2 \subseteq \cdots$ and taking a diagonal
subsequence gives uniform convergence on all compact subsets; the limit is holomorphic by Corollary~\ref{complex-analysis:corollary-liouville}(4).
\end{proof}

\begin{theorem}[Hurwitz]\label{complex-analysis:theorem-hurwitz}
Let $f_n \to f$ uniformly on compact subsets of a domain $U$, with $f_n$ holomorphic and $f$ nonconstant. If every $f_n$ is nowhere zero, so is $f$; if every $f_n$ is injective, so is $f$.
\end{theorem}

\begin{proof}
Suppose $f(z_0) = 0$. Since $f \not\equiv 0$, there is $r$ with $\overline{D}(z_0, r) \subseteq U$ and $f \neq 0$ on the circle $|z - z_0| = r$, so $|f| \geq \delta > 0$ there. Then $f_n \to f$ and $f_n' \to f'$
uniformly on the circle, so $\frac{1}{2\pi i}\int f_n'/f_n \to \frac{1}{2\pi i}\int f'/f \geq 1$, and by Theorem~\ref{complex-analysis:theorem-argument-principle} $f_n$ has a zero inside for large $n$. For
injectivity, if $f(z_1) = f(z_2)$ with $z_1 \neq z_2$, apply the first part on a punctured neighbourhood of $z_2$ to the functions $f_n - f_n(z_1)$, which are nowhere zero there.
\end{proof}

\begin{theorem}[Riemann mapping theorem]\label{complex-analysis:theorem-riemann-mapping}
Every simply connected domain $U \subsetneq \CC$ is biholomorphic to $\mathbb{D}$. More precisely, for $z_0 \in U$ there is a unique biholomorphic $\psi \colon U \to \mathbb{D}$ with $\psi(z_0) = 0$ and
$\psi'(z_0) > 0$.
\end{theorem}

\begin{proof}
Let $\mathcal{F}$ be the set of injective holomorphic maps $U \to \mathbb{D}$ sending $z_0$ to $0$.

\emph{$\mathcal{F} \neq \emptyset$.} Choose $b \notin U$. Since $z - b$ is nowhere zero on the simply connected $U$, the function $(z - b)'/(z - b)$ has a primitive $h$
(Theorem~\ref{complex-analysis:theorem-cauchy-homotopy}); adjusting by a constant, $e^{h} = z - b$. Let $\phi = e^{h/2}$, so $\phi^2 = z - b$. Then $\phi$ is injective, and $\phi(z_1) = -\phi(z_2)$ implies
$z_1 = z_2$ and $\phi(z_1) = 0$, which is impossible; so $\phi(U)$ and $-\phi(U)$ are disjoint. By Theorem~\ref{complex-analysis:theorem-open-mapping}, $\phi(U)$ contains a disc $D(w_0, \varrho)$, hence
$|\phi(z) + w_0| \geq \varrho$ on $U$, and $z \mapsto \varrho/(\phi(z) + w_0)$ maps $U$ injectively into $\mathbb{D}$. Composing with a suitable $\phi_a$
(Proposition~\ref{complex-analysis:proposition-disc-automorphisms}) gives an element of $\mathcal{F}$.

\emph{An extremal element.} Injective holomorphic maps have nonvanishing derivative, so $\eta = \sup_{\psi \in \mathcal{F}}|\psi'(z_0)| > 0$. Take $\psi_n \in \mathcal{F}$ with $|\psi_n'(z_0)| \to \eta$; by
Theorem~\ref{complex-analysis:theorem-montel} a subsequence converges uniformly on compact subsets to a holomorphic $\psi$ with $|\psi'(z_0)| = \eta > 0$, so $\psi$ is nonconstant, hence injective by
Theorem~\ref{complex-analysis:theorem-hurwitz}, with $\psi(z_0) = 0$ and $|\psi| \leq 1$; by the open mapping theorem $\psi(U) \subseteq \mathbb{D}$. So $\psi \in \mathcal{F}$ and $|\psi'(z_0)| = \eta$.

\emph{Surjectivity.} Suppose $\alpha \in \mathbb{D} \setminus \psi(U)$. Then $\phi_\alpha \circ \psi$ is nowhere zero on $U$, so as above it has a holomorphic square root $g$, which is injective because
$g^2$ is. Let $F = \phi_{g(z_0)} \circ g \in \mathcal{F}$. Writing $s(w) = w^2$, we have $\psi = \phi_{-\alpha} \circ s \circ \phi_{-g(z_0)} \circ F = \Phi \circ F$ with $\Phi \colon \mathbb{D} \to \mathbb{D}$
holomorphic and $\Phi(0) = 0$. As $s$ is not injective, $\Phi$ is not injective, so $\Phi$ is not a rotation and $|\Phi'(0)| < 1$ by Lemma~\ref{complex-analysis:lemma-schwarz}. Then
$\eta = |\psi'(z_0)| = |\Phi'(0)|\,|F'(z_0)| < |F'(z_0)| \leq \eta$, a contradiction. So $\psi$ is a biholomorphism onto $\mathbb{D}$, and multiplying by a unimodular constant makes $\psi'(z_0) > 0$.

\emph{Uniqueness.} If $\psi_1, \psi_2$ both work, $\psi_2 \circ \psi_1^{-1}$ is an automorphism of $\mathbb{D}$ fixing $0$ with positive derivative there, hence the identity by
Proposition~\ref{complex-analysis:proposition-disc-automorphisms}.
\end{proof}

\begin{remark}\label{complex-analysis:remark-harmonic}
A real function $u$ on an open $U \subseteq \CC$ is \emph{harmonic} if it is $C^2$ with $\partial_x^2u + \partial_y^2u = 0$. Real and imaginary parts of holomorphic functions are harmonic
(Corollary~\ref{complex-analysis:corollary-smooth-harmonic}). Conversely, on a simply connected $U$ every harmonic $u$ is the real part of a holomorphic function: the function
$\partial_xu - i\partial_yu$ is holomorphic (its Cauchy--Riemann equations are the harmonicity of $u$ and the symmetry of second derivatives), so it has a primitive $f$, and $\operatorname{Re}f$ differs from
$u$ by a constant. Harmonic functions therefore satisfy the mean value property and the maximum principle; this is the one-variable case of the Hodge theory of Chapter~\ref{hodge-theorem}.
\end{remark}

\section{The Riemann sphere}\label{complex-analysis:section-riemann-sphere}

Many statements of complex analysis become simpler once the point $\infty$ is added to the plane: a pole is a point where a function takes the value $\infty$, lines are circles through $\infty$,
and the rational functions are exactly the functions that are meromorphic everywhere, $\infty$ included. The resulting space is a sphere, the first example of a compact Riemann surface; its
symmetries, the M\"obius transformations, form the group $\mathrm{PSL}_2(\CC)$, which reappears in hyperbolic geometry, in the theory of modular forms and in the classification of Riemann surfaces.

\subsection{The sphere and stereographic projection}\label{complex-analysis:subsection-sphere}

\begin{definition}\label{complex-analysis:definition-riemann-sphere}
The \emph{Riemann sphere} is the one-point compactification $\hat{\CC} = \CC \cup \{\infty\}$ of $\CC$ (Definition~\ref{general-topology:definition-one-point-compactification}). The neighbourhoods of $\infty$ are the sets
containing $\{|z| > R\} \cup \{\infty\}$ for some $R$, so $z_n \to \infty$ in $\hat{\CC}$ if and only if $|z_n| \to \infty$.
\end{definition}

Identifying $\CC$ with $\RR^2$ by $z = x + iy \mapsto (x, y)$, Example~\ref{general-topology:example-stereographic} gives a homeomorphism $\tau \colon \hat{\CC} \to S^2 \subseteq \RR^3$,
\[
\tau(z) = \frac{(2x, 2y, |z|^2 - 1)}{|z|^2 + 1} = \frac{(2\operatorname{Re}z, 2\operatorname{Im}z, |z|^2 - 1)}{|z|^2 + 1}, \qquad \tau(\infty) = N = (0, 0, 1),
\]
whose inverse is the stereographic projection from the north pole $N$: the point $\tau(z)$ is the second intersection of $S^2$ with the line through $N$ and $(x, y, 0)$. The unit circle is mapped to the
equator, the unit disc $\mathbb{D}$ to the southern hemisphere, $0$ to the south pole and $\infty$ to the north pole.

\begin{proposition}\label{complex-analysis:proposition-chordal}
\begin{enumerate}
\item (\emph{Chordal distance}) For $z, w \in \CC$,
\[
\|\tau(z) - \tau(w)\| = \frac{2|z - w|}{\sqrt{(1 + |z|^2)(1 + |w|^2)}}, \qquad \|\tau(z) - N\| = \frac{2}{\sqrt{1 + |z|^2}}.
\]
\item The map $z \mapsto 1/z$, with $0 \mapsto \infty \mapsto 0$, corresponds to the rotation by the angle $\pi$ about the $x_1$-axis: $\tau(1/z) = (x_1, -x_2, -x_3)$ if $\tau(z) = (x_1, x_2, x_3)$. In particular it
preserves the chordal distance.
\end{enumerate}
\end{proposition}

\begin{proof}
(1) Since $\|\tau(z)\| = \|\tau(w)\| = 1$, $\|\tau(z) - \tau(w)\|^2 = 2 - 2\langle\tau(z), \tau(w)\rangle$, and
\[
2 - 2\langle\tau(z), \tau(w)\rangle = 2\,\frac{(|z|^2 + 1)(|w|^2 + 1) - 4\operatorname{Re}(z\bar w) - (|z|^2 - 1)(|w|^2 - 1)}{(|z|^2 + 1)(|w|^2 + 1)}
= \frac{4(|z|^2 + |w|^2 - 2\operatorname{Re}(z\bar w))}{(|z|^2 + 1)(|w|^2 + 1)},
\]
which is $4|z - w|^2/((1 + |z|^2)(1 + |w|^2))$. The second formula is in Example~\ref{general-topology:example-stereographic}.

(2) For $z \neq 0$, $1/z = \bar z/|z|^2$ has real part $x/|z|^2$ and imaginary part $-y/|z|^2$, so
\[
\tau(1/z) = \frac{(2x/|z|^2, -2y/|z|^2, 1/|z|^2 - 1)}{1/|z|^2 + 1} = \frac{(2x, -2y, 1 - |z|^2)}{1 + |z|^2},
\]
which is $(x_1, -x_2, -x_3)$; and $\tau$ maps $0$ and $\infty$ to the south and north poles, which the rotation exchanges.
\end{proof}

A \emph{generalised circle} in $\hat{\CC}$ is either a circle in $\CC$ or a set $L \cup \{\infty\}$ with $L$ a line in $\CC$. Every generalised circle is the set of solutions in $\hat{\CC}$ of an equation
\begin{equation}\label{complex-analysis:equation-generalised-circle}
A|z|^2 + 2\beta x + 2\gamma y + C = 0, \qquad A, \beta, \gamma, C \in \RR, \quad \beta^2 + \gamma^2 > AC,
\end{equation}
where $\infty$ counts as a solution when $A = 0$: for $A \neq 0$ this is the circle with centre $-(\beta + i\gamma)/A$ and radius $\sqrt{\beta^2 + \gamma^2 - AC}/|A|$, and for $A = 0$ a line. Through any three
distinct points of $\hat{\CC}$ passes exactly one generalised circle: the circle through three non-collinear points of $\CC$, the line through three collinear ones, and the line through the two finite
points if one of them is $\infty$.

\begin{proposition}\label{complex-analysis:proposition-stereographic-circles}
The map $\tau$ induces a bijection between the generalised circles in $\hat{\CC}$ and the circles on $S^2$, that is, the intersections of $S^2$ with planes at distance less than $1$ from the origin. The lines
correspond to the circles through $N$.
\end{proposition}

\begin{proof}
A plane is $\{a x_1 + bx_2 + cx_3 = d\}$ with $a^2 + b^2 + c^2 = 1$, and it meets $S^2$ in a circle exactly when $|d| < 1$. The point $\tau(z)$, $z \in \CC$, lies on it if and only if
$2ax + 2by + c(|z|^2 - 1) = d(|z|^2 + 1)$, that is,
\[
(c - d)|z|^2 + 2ax + 2by - (c + d) = 0,
\]
an equation of the form~\eqref{complex-analysis:equation-generalised-circle} with $A = c - d$, $\beta = a$, $\gamma = b$, $C = -(c + d)$, for which $\beta^2 + \gamma^2 - AC = a^2 + b^2 + c^2 - d^2 = 1 - d^2 > 0$. The
point $N = \tau(\infty)$ lies on the plane if and only if $c = d$, that is, $A = 0$. Conversely, given $(A, \beta, \gamma, C)$ with $\beta^2 + \gamma^2 > AC$, the numbers $a = s\beta$, $b = s\gamma$, $c = s(A - C)/2$,
$d = -s(A + C)/2$ satisfy the equations above, and $a^2 + b^2 + c^2 - d^2 = s^2(\beta^2 + \gamma^2 - AC) > 0$; choosing $s > 0$ with $a^2 + b^2 + c^2 = 1$ gives a plane with $1 - d^2 > 0$ whose circle corresponds to the
given generalised circle.
\end{proof}

Stereographic projection also preserves angles between curves; equivalently, the round metric of $S^2$ pulls back to the conformal metric $4|dz|^2/(1 + |z|^2)^2$ (Chapter~\ref{riemannian-geometry}).

\subsection{Meromorphic functions on the sphere}\label{complex-analysis:subsection-rational}

The behaviour of a function at $\infty$ is defined through the coordinate $w = 1/z$.

\begin{definition}\label{complex-analysis:definition-infinity}
Let $f$ be holomorphic on $\{|z| > R\}$, and $g(w) = f(1/w)$ on $0 < |w| < 1/R$. Then $f$ is \emph{holomorphic at $\infty$}, has a \emph{pole}, a \emph{zero} or an \emph{essential singularity} at
$\infty$, if $g$ has the corresponding property at $0$, and $\operatorname{ord}_\infty f = \operatorname{ord}_0g$. In terms of the Laurent expansion $f(z) = \sum_{n \in \ZZ}c_nz^n$ on $|z| > R$
(Theorem~\ref{complex-analysis:theorem-laurent}), where $g(w) = \sum_nc_nw^{-n}$: $f$ is holomorphic at $\infty$ if $c_n = 0$ for $n > 0$, and has a pole of order $m$ if $c_m \neq 0$ and $c_n = 0$ for $n > m$. The
\emph{residue at $\infty$} is
\[
\operatorname{Res}_\infty f = -\frac{1}{2\pi i}\int_{|z| = \rho}f(z)\,dz = -c_{-1} \qquad (\rho > R),
\]
the circle traversed counterclockwise. A function $f$ on $\hat{\CC}$ is \emph{meromorphic} if it is meromorphic on $\CC$ and holomorphic or has a pole at $\infty$; it is regarded as a map $\hat{\CC} \to \hat{\CC}$ by sending
the poles to $\infty$, which is continuous by Proposition~\ref{complex-analysis:proposition-singularities}(2).
\end{definition}

The residue at $\infty$ is not the residue of $g$ at $0$: it belongs to the differential $f(z)\,dz = -f(1/w)w^{-2}\,dw$, and $\operatorname{Res}_\infty f = \operatorname{Res}_0(-g(w)/w^2)$. For instance $f(z) = 1/z$ is
holomorphic at $\infty$, with a zero of order $1$ there, but $\operatorname{Res}_\infty(1/z) = -1$. The minus sign is there because the circle $|z| = \rho$, traversed counterclockwise, goes around $\infty$ clockwise.

A rational function $R = p/q$, with $p, q$ coprime polynomials and $q \neq 0$, is meromorphic on $\hat{\CC}$: near $\infty$, $R(z) = z^{\deg p - \deg q}h(1/z)$ with $h$ holomorphic near $0$ and $h(0) \neq 0$ (the ratio
of the leading coefficients), so
\begin{equation}\label{complex-analysis:equation-order-infinity}
\operatorname{ord}_\infty\frac pq = \deg q - \deg p.
\end{equation}
The \emph{degree} of $R$ is $\deg R = \max(\deg p, \deg q)$.

\begin{theorem}\label{complex-analysis:theorem-rational}
The meromorphic functions on $\hat{\CC}$ are exactly the rational functions. More precisely, if $f$ is meromorphic on $\hat{\CC}$, with poles $a_1, \ldots, a_k$ in $\CC$, principal parts $P_j(z) = \sum_{n=1}^{m_j}c^{(j)}_{-n}(z - a_j)^{-n}$ there,
and principal part $P_\infty(z) = \sum_{n=1}^mc_nz^n$ at $\infty$ (the positive part of its Laurent expansion at $\infty$), then
\[
f = c + P_\infty + P_1 + \cdots + P_k
\]
for a constant $c$. This is the \emph{partial fraction decomposition} of $f$, and it is unique.
\end{theorem}

\begin{proof}
Let $f = \sum_nc_nz^n$ on $|z| > R$. Since $f$ is not essential at $\infty$, $c_n = 0$ for large $n$, so $P_\infty$ is a polynomial. The poles of $f$ in $\CC$ lie in the compact disc $|z| \leq R$, where $f$ is
holomorphic outside a discrete closed set of poles, so there are finitely many, $a_1, \ldots, a_k$. Each $P_j$ is a rational function, holomorphic on $\hat{\CC} \setminus \{a_j\}$ and vanishing at $\infty$. The function
$h = f - P_\infty - \sum_jP_j$ has removable singularities at the $a_j$, so it extends to an entire function. As $z \to \infty$, $f(z) - P_\infty(z) = \sum_{n \leq 0}c_nz^n \to c_0$ and $P_j(z) \to 0$, so $h$ is bounded
outside a large disc, hence bounded on $\CC$, and constant by Liouville's theorem (Corollary~\ref{complex-analysis:corollary-liouville}). So $f$ is rational. Uniqueness: the principal parts, and hence $c$,
are determined by $f$.
\end{proof}

So a meromorphic function on $\hat{\CC}$ is determined, up to an additive constant, by its principal parts, and up to a multiplicative constant by its zeros and poles: if $f$ and $g$ have the same zeros and
poles with the same orders, $f/g$ is meromorphic on $\hat{\CC}$ with only removable singularities, so it has no principal parts and is constant by the theorem. This is the simplest instance of the
interplay of local data and global holomorphic functions studied for general compact Riemann surfaces in Chapter~\ref{curves}.

\begin{proposition}\label{complex-analysis:proposition-rational-degree}
Let $R$ be a nonconstant rational function of degree $d$. For every $w \in \hat{\CC}$, the equation $R(z) = w$ has exactly $d$ solutions in $\hat{\CC}$, counted with multiplicity; here the multiplicity of a solution
$a$ is $\operatorname{ord}_a(R - w)$ if $a$ and $w$ are finite, and is defined through the coordinates $1/z$ near $z = \infty$ and $1/R$ near $w = \infty$ otherwise. Consequently
$\sum_{a \in \hat{\CC}}\operatorname{ord}_aR = 0$: a rational function has as many zeros as poles on $\hat{\CC}$, counted with multiplicity.
\end{proposition}

\begin{proof}
Write $R = p/q$ with $p$, $q$ coprime. Let $w \in \CC$. The finite solutions of $R(z) = w$ are the roots of $p - wq$: at a root $a$ of $p - wq$ we have $q(a) \neq 0$, as otherwise $p(a) = 0$ too, and then
$\operatorname{ord}_a(R - w) = \operatorname{ord}_a(p - wq)$. By the fundamental theorem of algebra, these multiplicities add up to $e = \deg(p - wq) \leq d$. The point $\infty$ is a solution exactly when
$R - w = (p - wq)/q$ vanishes at $\infty$, with multiplicity $\deg q - e$ by~\eqref{complex-analysis:equation-order-infinity}. If $\deg p > \deg q$, then $e = \deg p = d$ and $R(\infty) = \infty \neq w$. If
$\deg p \leq \deg q = d$, then either $e = d$, and $R - w$ does not vanish at $\infty$; or $e < d$, and $\infty$ is a solution of multiplicity $d - e$. In all cases the total is $d$. For $w = \infty$, the solutions
are the poles: the roots of $q$, with multiplicities adding up to $\deg q$, and $\infty$ if $\deg p > \deg q$, with multiplicity $\deg p - \deg q$; the total is again $d$. Finally $\sum_a\operatorname{ord}_aR$ is the number
of zeros minus the number of poles, $d - d = 0$.
\end{proof}

\begin{proposition}\label{complex-analysis:proposition-residues-sphere}
For every rational function $R$, $\sum_{a \in \hat{\CC}}\operatorname{Res}_aR = 0$, the sum running over $\infty$ and the poles of $R$ in $\CC$.
\end{proposition}

\begin{proof}
Choose $\rho$ larger than the absolute values of all poles of $R$ in $\CC$. By the residue theorem in the simply connected set $\CC$ (Theorem~\ref{complex-analysis:theorem-residues}),
$\frac{1}{2\pi i}\int_{|z| = \rho}R(z)\,dz = \sum_{a \in \CC}\operatorname{Res}_aR$, and the left side is $-\operatorname{Res}_\infty R$ by definition.
\end{proof}

\begin{example}\label{complex-analysis:example-residue-infinity}
\begin{enumerate}
\item The residue at $\infty$ often computes a contour integral at once. For $n \geq 2$, $\int_{|z| = 2}\frac{dz}{z^n - 1} = 0$, because $\frac{1}{z^n - 1} = z^{-n}(1 + z^{-n} + z^{-2n} + \cdots)$ for $|z| > 1$ has
$c_{-1} = 0$; computing the $n$ residues at the roots of unity would be longer. Similarly $\int_{|z| = 2}\frac{z^{n-1}\,dz}{z^n - 1} = 2\pi i$, which is also $2\pi i$ times the number of roots of $z^n - 1$ in the
disc, by the argument principle (Theorem~\ref{complex-analysis:theorem-argument-principle}).
\item Applying the proposition to $R(z) = \frac{1}{z(z - 1)}$: the residues at $0$ and $1$ are $-1$ and $1$, and indeed $R = z^{-2}(1 + z^{-1} + \cdots)$ has $\operatorname{Res}_\infty R = 0$.
\item A polynomial of degree $d \geq 1$ is a rational function of degree $d$ with a single pole, of order $d$, at $\infty$; Proposition~\ref{complex-analysis:proposition-rational-degree} says it takes every value $d$
times in $\CC$, which is the fundamental theorem of algebra.
\end{enumerate}
\end{example}

\subsection{M\"obius transformations}\label{complex-analysis:subsection-mobius}

\begin{definition}\label{complex-analysis:definition-mobius}
For $A = \bigl(\begin{smallmatrix}a & b\\ c & d\end{smallmatrix}\bigr) \in \GL_2(\CC)$, the \emph{M\"obius transformation} $M_A \colon \hat{\CC} \to \hat{\CC}$ is
\[
M_A(z) = \frac{az + b}{cz + d},
\]
with $M_A(-d/c) = \infty$ and $M_A(\infty) = a/c$ if $c \neq 0$, and $M_A(\infty) = \infty$ if $c = 0$.
\end{definition}

It is convenient to write points of $\hat{\CC}$ in \emph{homogeneous coordinates}: the nonzero vector $(z_1, z_2) \in \CC^2$ represents $z_1/z_2 \in \hat{\CC}$ (with $z_1/0 = \infty$), and two vectors represent the same
point if and only if they are proportional. So $\hat{\CC}$ is the set of lines through $0$ in $\CC^2$, the \emph{complex projective line} $\PP^1(\CC)$, and $M_A$ is the map induced by the linear map $A$: if
$(z_1, z_2)$ represents $z$, then $A(z_1, z_2) = (az_1 + bz_2, cz_1 + dz_2)$ represents $M_A(z)$. (The cases with $\infty$ in the definition are exactly what this description gives.)

\begin{proposition}\label{complex-analysis:proposition-mobius}
\begin{enumerate}
\item $M_AM_B = M_{AB}$, and $M_A = \id$ if and only if $A$ is a scalar matrix. So the M\"obius transformations form a group, isomorphic to $\mathrm{PGL}_2(\CC) = \GL_2(\CC)/\CC^*$, and also to
$\mathrm{PSL}_2(\CC) = \SL_2(\CC)/\{\pm1\}$.
\item Each $M_A$ is a homeomorphism of $\hat{\CC}$ and a rational function of degree $1$; its derivative $\det A/(cz + d)^2$ never vanishes on $\CC \setminus \{-d/c\}$. Conversely every rational function of degree $1$, and
every bijective meromorphic function $\hat{\CC} \to \hat{\CC}$, is a M\"obius transformation.
\item The M\"obius transformations are generated by the translations $z \mapsto z + b$, the maps $z \mapsto az$ ($a \neq 0$) and the inversion $z \mapsto 1/z$.
\item M\"obius transformations map generalised circles to generalised circles.
\end{enumerate}
\end{proposition}

\begin{proof}
(1) The first statement is the description in homogeneous coordinates: $AB$ induces the composite of the maps induced by $B$ and $A$. If $M_A = \id$, every line in $\CC^2$ is mapped to itself, so every
nonzero vector is an eigenvector of $A$; then $e_1$, $e_2$ and $e_1 + e_2$ are eigenvectors, which forces the two eigenvalues for $e_1$, $e_2$ to agree, and $A$ is scalar. So $A \mapsto M_A$ induces an isomorphism
from $\GL_2(\CC)/\CC^*$ onto the group of M\"obius transformations. Every $A$ is a scalar multiple of a matrix of determinant $1$ (divide by a square root of $\det A$), and a scalar matrix of determinant $1$ is
$\pm1$, which gives the second description.

(2) $M_A$ is bijective with inverse $M_{A^{-1}}$. It is continuous on $\hat{\CC}$: away from $-d/c$ and $\infty$ this is clear, $|M_A(z)| \to \infty$ as $z \to -d/c$ since the numerator does not vanish there
($a(-d/c) + b = -\det A/c$), and $M_A(z) \to a/c$ (or $\infty$ if $c = 0$) as $z \to \infty$. A continuous bijection of the compact Hausdorff space $\hat{\CC}$ is a homeomorphism. The polynomials $az + b$ and $cz + d$
are coprime, since $ad - bc \neq 0$, so $M_A$ has degree $1$, and the formula for the derivative is the quotient rule. Conversely, a rational function $(az + b)/(cz + d)$ of degree $1$ in lowest terms has
$ad - bc \neq 0$, since otherwise the numerator and denominator would be proportional. If $f$ is a meromorphic bijection of $\hat{\CC}$, it is rational (Theorem~\ref{complex-analysis:theorem-rational}) of some degree
$d \geq 1$, and by Proposition~\ref{complex-analysis:proposition-rational-degree} each value has a single preimage, of multiplicity $d$. If $d \geq 2$, then $f^{\prime}$ would vanish at every point of $\CC$ that is not a pole and
not mapped to $\infty$, which is impossible for a nonconstant function. So $d = 1$.

(3) If $c = 0$, $M_A(z) = (a/d)z + b/d$. If $c \neq 0$, $M_A(z) = \frac ac - \frac{\det A}{c(cz + d)}$, the composite of $z \mapsto cz + d$, $z \mapsto 1/z$, multiplication by $-\det A/c$ and translation by $a/c$.

(4) By (3) it suffices to check the generators. Translations and the maps $z \mapsto az$ are similarities of the plane, which map circles to circles and lines to lines, and fix $\infty$. The inversion corresponds under
$\tau$ to a rotation of $S^2$ (Proposition~\ref{complex-analysis:proposition-chordal}(2)), which maps plane sections of $S^2$ to plane sections; by Proposition~\ref{complex-analysis:proposition-stereographic-circles}
it maps generalised circles to generalised circles.
\end{proof}

The proof of (4) by the sphere avoids a computation; directly, substituting $z = 1/u$ in~\eqref{complex-analysis:equation-generalised-circle} and multiplying by $|u|^2$ gives
$A + 2\beta\operatorname{Re}u - 2\gamma\operatorname{Im}u + C|u|^2 = 0$, again of the same form. Under inversion, circles through $0$ become lines, and lines not through $0$ become circles through $0$.

\begin{theorem}\label{complex-analysis:theorem-three-points}
Given distinct points $z_1, z_2, z_3 \in \hat{\CC}$ and distinct points $w_1, w_2, w_3 \in \hat{\CC}$, there is exactly one M\"obius transformation $M$ with $M(z_i) = w_i$ for $i = 1, 2, 3$.
\end{theorem}

\begin{proof}
If $z_1, z_2, z_3 \in \CC$, the map
\[
S(z) = \frac{(z - z_1)(z_2 - z_3)}{(z - z_3)(z_2 - z_1)}
\]
is M\"obius, with $S(z_1) = 0$, $S(z_2) = 1$, $S(z_3) = \infty$. If one of the $z_i$ is $\infty$, delete the two factors containing it: $S(z) = \frac{z_2 - z_3}{z - z_3}$, $\frac{z - z_1}{z - z_3}$ or $\frac{z - z_1}{z_2 - z_1}$ for
$z_1$, $z_2$ or $z_3 = \infty$. Doing the same for the $w_i$ gives $T$, and $M = T^{-1}S$ has the required property. For uniqueness it suffices to show that a M\"obius map fixing $0$, $1$, $\infty$ is the identity: from
$M(\infty) = \infty$, $c = 0$, so $M(z) = \alpha z + \beta$; then $M(0) = 0$ gives $\beta = 0$ and $M(1) = 1$ gives $\alpha = 1$.
\end{proof}

\begin{definition}\label{complex-analysis:definition-cross-ratio}
For distinct $z_1, z_2, z_3 \in \hat{\CC}$ and $z \in \hat{\CC}$, the \emph{cross-ratio} $(z, z_1, z_2, z_3)$ is $S(z)$, for the unique M\"obius $S$ with $S(z_1) = 0$, $S(z_2) = 1$, $S(z_3) = \infty$; for finite points
\[
(z, z_1, z_2, z_3) = \frac{(z - z_1)(z_2 - z_3)}{(z - z_3)(z_2 - z_1)}.
\]
\end{definition}

\begin{corollary}\label{complex-analysis:corollary-cross-ratio}
\begin{enumerate}
\item For every M\"obius transformation $M$, $(Mz, Mz_1, Mz_2, Mz_3) = (z, z_1, z_2, z_3)$.
\item Four distinct points $z, z_1, z_2, z_3$ lie on a generalised circle if and only if $(z, z_1, z_2, z_3)$ is real.
\end{enumerate}
\end{corollary}

\begin{proof}
(1) $SM^{-1}$ sends $Mz_1$, $Mz_2$, $Mz_3$ to $0$, $1$, $\infty$, so it is the map defining the left side, and its value at $Mz$ is $S(z)$. (2) Let $C$ be the generalised circle through $z_1$, $z_2$, $z_3$. By
Proposition~\ref{complex-analysis:proposition-mobius}(4), applied to $S$ and $S^{-1}$, $S$ maps $C$ onto the generalised circle through $0$, $1$, $\infty$, which is $\RR \cup \{\infty\}$. So $z \in C$ if and only if
$S(z) \in \RR \cup \{\infty\}$, and $S(z) \neq \infty$ since $z \neq z_3$.
\end{proof}

\subsection{Fixed points and conjugacy classes}\label{complex-analysis:subsection-mobius-classification}

A point $z = z_1/z_2$ is fixed by $M_A$ if and only if $(z_1, z_2)$ is an eigenvector of $A$. Since $A$ has one or two eigenlines unless it is scalar, a M\"obius transformation other than the identity has one or two
fixed points in $\hat{\CC}$. Normalising $\det A = 1$, the trace $\operatorname{tr}A$ is determined by $M_A$ up to sign, so $(\operatorname{tr}A)^2$ is an invariant of $M_A$, unchanged under
conjugation $M \mapsto NMN^{-1}$.

\begin{theorem}\label{complex-analysis:theorem-mobius-classification}
Let $M = M_A \neq \id$ with $\det A = 1$, and $t = \operatorname{tr}A$.
\begin{enumerate}
\item If $t^2 = 4$, $M$ has exactly one fixed point and is conjugate to the translation $z \mapsto z + 1$; $M$ is called \emph{parabolic}.
\item If $t^2 \neq 4$, $M$ has exactly two fixed points and is conjugate to $z \mapsto kz$ for some $k \neq 0, 1$, determined up to $k \leftrightarrow k^{-1}$ by $k + 2 + k^{-1} = t^2$. Then $M$ is called \emph{elliptic} if
$|k| = 1$, equivalently $t \in (-2, 2)$; \emph{hyperbolic} if $k > 0$, equivalently $t \in \RR$ with $|t| > 2$; and \emph{loxodromic} otherwise, equivalently $t \notin \RR$.
\end{enumerate}
Two M\"obius transformations other than the identity are conjugate if and only if they have the same invariant $t^2$.
\end{theorem}

\begin{proof}
Conjugating $M_A$ by $M_P$ replaces $A$ by $PAP^{-1}$. By the Jordan normal form (Theorem~\ref{linear-algebra:theorem-jordan}), a non-scalar $A \in \SL_2(\CC)$ is conjugate either to
$\bigl(\begin{smallmatrix}\lambda & 0\\ 0 & \lambda^{-1}\end{smallmatrix}\bigr)$ with $\lambda \neq \lambda^{-1}$, or to $\bigl(\begin{smallmatrix}\lambda & 1\\ 0 & \lambda\end{smallmatrix}\bigr)$ with $\lambda^2 = 1$; a diagonal matrix
with $\lambda = \lambda^{-1}$ would be scalar. In the second case $t = 2\lambda = \pm2$, $M$ is conjugate to $z \mapsto z \pm 1$, and $z \mapsto z - 1$ is conjugate to $z \mapsto z + 1$ by $z \mapsto -z$; the only
fixed point of a translation is $\infty$. In the first case $M$ is conjugate to $z \mapsto kz$ with $k = \lambda^2 \neq 1$, which fixes $0$ and $\infty$, and $t^2 = (\lambda + \lambda^{-1})^2 = k + 2 + k^{-1}$; here $t^2 \neq 4$
since $k \neq 1$. Exchanging the eigenvectors replaces $k$ by $k^{-1}$. The eigenvalues $\lambda$, $\lambda^{-1}$ are the roots of $X^2 - tX + 1$, and $k = \lambda^2$. If $t \in (-2, 2)$, the roots are non-real and
conjugate with product $1$, so $|\lambda| = 1$ and $|k| = 1$; conversely $|\lambda| = 1$ gives $t = \lambda + \bar\lambda = 2\operatorname{Re}\lambda \in [-2, 2]$, and $t \neq \pm2$ as $t^2 \neq 4$. If $t$ is real with $|t| > 2$, the
roots are real, so $k > 0$; conversely $k > 0$ means $\lambda \in \RR$, so $t = \lambda + \lambda^{-1}$ is real, and $|t| > 2$ since $\lambda \neq \pm1$. As $t^2 \neq 4$, a real $t$ falls in one of these two cases, so the
remaining, loxodromic, case is $t \notin \RR$. Finally, transformations with the same $t^2$ are
conjugate, since both are conjugate to the same normal form ($z \mapsto z + 1$, or $z \mapsto kz$ with $k$ determined up to inversion, and $z \mapsto k^{-1}z$ is conjugate to $z \mapsto kz$ by $z \mapsto 1/z$).
\end{proof}

The names describe the dynamics. A hyperbolic map $z \mapsto kz$, $k > 1$, pushes every point away from the repelling fixed point $0$ towards the attracting fixed point $\infty$ along rays; an elliptic map
$z \mapsto e^{i\theta}z$ rotates around its two fixed points, and corresponds to a rotation of the sphere about the axis through them; a loxodromic map spirals from one fixed point to the other; and a parabolic map
moves points along circles tangent at the single fixed point. For example $z \mapsto 1/z$, with matrix $i\bigl(\begin{smallmatrix}0 & 1\\ 1 & 0\end{smallmatrix}\bigr)$ of determinant $1$ and trace $0$, is elliptic of
order $2$ with fixed points $\pm1$, in accordance with Proposition~\ref{complex-analysis:proposition-chordal}(2).

\subsection{The disc, the half-plane and the sphere}\label{complex-analysis:subsection-mobius-subgroups}

Three subgroups of the M\"obius group correspond to the three simply connected Riemann surfaces of the uniformisation theorem (Chapter~\ref{complex-manifolds}): the sphere, the plane and the disc.

\begin{proposition}\label{complex-analysis:proposition-half-plane}
Let $\mathbb{H} = \{\operatorname{Im}z > 0\}$.
\begin{enumerate}
\item The \emph{Cayley transform} $\kappa(z) = \frac{z - i}{z + i}$ maps $\mathbb{H}$ biholomorphically onto $\mathbb{D}$, and $\RR \cup \{\infty\}$ onto the unit circle.
\item For $A \in \SL_2(\RR)$, $\operatorname{Im}M_A(z) = \operatorname{Im}z/|cz + d|^2$, so $M_A$ maps $\mathbb{H}$ onto itself.
\item Every biholomorphic self-map of $\mathbb{H}$ is $M_A$ for some $A \in \SL_2(\RR)$, unique up to sign; so $\mathrm{Aut}(\mathbb{H}) \cong \mathrm{PSL}_2(\RR)$.
\end{enumerate}
\end{proposition}

\begin{proof}
(1) $|\kappa(z)| < 1$ if and only if $|z - i| < |z + i|$, that is, $z$ is closer to $i$ than to $-i$, that is, $\operatorname{Im}z > 0$; similarly $|\kappa(z)| = 1$ exactly on $\RR$, and $\kappa(\infty) = 1$. Since $\kappa$ is a
bijection of $\hat{\CC}$, it maps $\mathbb{H}$ onto $\mathbb{D}$, and it is holomorphic with holomorphic inverse $M_{A^{-1}}$.

(2) $(az + b)(c\bar z + d) = ac|z|^2 + adz + bc\bar z + bd$ has imaginary part $(ad - bc)\operatorname{Im}z = \operatorname{Im}z$, and $\operatorname{Im}M_A(z) = \operatorname{Im}\bigl((az + b)(c\bar z + d)\bigr)/|cz + d|^2$. So $M_A$
maps $\mathbb{H}$ into itself, and so does its inverse $M_{A^{-1}}$, $A^{-1} \in \SL_2(\RR)$.

(3) Let $f$ be a biholomorphic self-map of $\mathbb{H}$. Then $\kappa f\kappa^{-1}$ is a biholomorphic self-map of $\mathbb{D}$, hence a M\"obius transformation by
Proposition~\ref{complex-analysis:proposition-disc-automorphisms}, and so $f = M$ is M\"obius. As a homeomorphism of $\hat{\CC}$ mapping $\mathbb{H}$ onto itself, $M$ maps the boundary $\RR \cup \{\infty\}$ of $\mathbb{H}$ onto
itself. Choose three distinct real points $z_i$; the $w_i = M(z_i)$ lie in $\RR \cup \{\infty\}$, and in the proof of Theorem~\ref{complex-analysis:theorem-three-points} the maps $S$ and $T$ have real coefficients, so
$M = T^{-1}S = M_A$ with $A \in \GL_2(\RR)$. By the computation in (2), $\operatorname{Im}M_A(z) = \det A\operatorname{Im}z/|cz + d|^2$, so $\det A > 0$, and dividing $A$ by $\sqrt{\det A}$ gives $A \in \SL_2(\RR)$. Finally
$M_A = M_B$ with $A, B \in \SL_2(\RR)$ forces $A = \pm B$ by Proposition~\ref{complex-analysis:proposition-mobius}(1).
\end{proof}

So the three groups are: all of $\mathrm{PSL}_2(\CC)$ for the sphere (Proposition~\ref{complex-analysis:proposition-mobius}(2)); the affine maps $z \mapsto az + b$, the M\"obius maps fixing $\infty$, for the plane
(Exercise~\ref{complex-analysis:exercise-automorphisms-plane}); and $\mathrm{PSL}_2(\RR) \cong \mathrm{PSU}(1, 1)$ for the half-plane or the disc. The subgroup $\mathrm{PSL}_2(\ZZ)$ acting on $\mathbb{H}$ gives the moduli
space of elliptic curves (Chapter~\ref{deformations}).

The rotations of the sphere form a fourth, compact, subgroup. In homogeneous coordinates the chordal distance of Proposition~\ref{complex-analysis:proposition-chordal} reads
\[
\|\tau(z) - \tau(w)\| = \frac{2|z_1w_2 - z_2w_1|}{\|(z_1, z_2)\|\,\|(w_1, w_2)\|} \qquad (z = z_1/z_2, \ w = w_1/w_2),
\]
as one sees by taking $z_2 = w_2 = 1$ (both sides are unchanged by rescaling either vector, and the case of $\infty$ follows by continuity). The numerator is $|\det|$ of the matrix with columns $(z_1, z_2)$ and
$(w_1, w_2)$, so a matrix $A \in \mathrm{SU}(2)$, which preserves determinants and norms, gives a M\"obius transformation $M_A$ preserving the chordal distance, an isometry of $S^2$. These isometries are rotations, and
every rotation arises in this way; this is the double covering $\mathrm{SU}(2) \to \mathrm{SO}(3)$ of Chapter~\ref{lie-groups}. The matrices in $\mathrm{SU}(2)$ are those of the form
$\bigl(\begin{smallmatrix}a & -\bar b\\ b & \bar a\end{smallmatrix}\bigr)$ with $|a|^2 + |b|^2 = 1$; for instance $z \mapsto e^{i\theta}z$ is the rotation by $\theta$ about the vertical axis.

\begin{remark}\label{complex-analysis:remark-riemann-sphere}
\begin{enumerate}
\item \emph{Riemann surfaces.} With the two charts $z$ and $1/z$, $\hat{\CC}$ is a compact Riemann surface, biholomorphic to the projective line $\CC\PP^1$; meromorphic functions on $\hat{\CC}$ are the holomorphic maps
$\hat{\CC} \to \hat{\CC}$ other than the constant $\infty$, and the M\"obius transformations are its automorphisms (Chapter~\ref{complex-manifolds}). It is the only compact Riemann surface of genus $0$.
\item \emph{Riemann--Hurwitz.} A rational function of degree $d$ has $2d - 2$ critical points on $\hat{\CC}$, counted with multiplicity; for a polynomial of degree $d$ these are the $d - 1$ zeros of $p^{\prime}$ and the point
$\infty$, of multiplicity $d - 1$. This is the Riemann--Hurwitz formula for maps $\hat{\CC} \to \hat{\CC}$ \cite{forster-riemann-surfaces}.
\item \emph{Iteration.} Iterating a rational function of degree $d \geq 2$ on $\hat{\CC}$ is the subject of complex dynamics, with its Julia and Fatou sets; the family $z \mapsto z^2 + c$ gives the Mandelbrot set.
For M\"obius transformations iteration is simple, by Theorem~\ref{complex-analysis:theorem-mobius-classification}: the iterates of $z \mapsto kz$ are $z \mapsto k^nz$, and those of $z \mapsto z + 1$ are $z \mapsto z + n$.
\end{enumerate}
\end{remark}

\section{Exercises}\label{complex-analysis:section-exercises}

\begin{exercise}\label{complex-analysis:exercise-disc-not-plane}
Show that $\mathbb{D}$ and $\CC$ are not biholomorphic, although they are homeomorphic.
\end{exercise}

\begin{solution}
A biholomorphism $\CC \to \mathbb{D}$ would be a bounded entire function, hence constant (Corollary~\ref{complex-analysis:corollary-liouville}). They are homeomorphic by
$z \mapsto z/(1 + |z|)$ (Counterexample~\ref{general-topology:counterexample-completeness-not-topological}).
\end{solution}

\begin{exercise}\label{complex-analysis:exercise-automorphisms-plane}
Show that every biholomorphic map $\CC \to \CC$ is $z \mapsto az + b$ with $a \neq 0$.
\end{exercise}

\begin{solution}
Let $f$ be such a map and consider $g(z) = f(1/z)$ on $\CC^\times$. If $0$ were an essential singularity of $g$, then by Casorati--Weierstrass the image $f(\{|z| > 1\})$ would be dense, while $f(\mathbb{D})$
is open and disjoint from it, a contradiction. So $g$ has a pole or removable singularity at $0$, which means $|f(z)| \leq C|z|^m$ for large $|z|$; by Liouville's theorem $f$ is a polynomial, and it is
injective only if $\deg f = 1$.
\end{solution}

\begin{exercise}\label{complex-analysis:exercise-rouche-roots}
Use Rouch\'e's theorem to show that $z^5 + 3z + 1$ has exactly one root in $\mathbb{D}$ and all five roots in $D(0, 2)$.
\end{exercise}

\begin{exercise}\label{complex-analysis:exercise-residue-trigonometric}
Compute $\int_0^{2\pi}\frac{dt}{a + \cos t}$ for $a > 1$ by writing $z = e^{it}$ and using the residue theorem.
\end{exercise}

\begin{exercise}\label{complex-analysis:exercise-cauchy-not-converse}
Give an example of a domain $U$ and a nowhere vanishing holomorphic $f$ on $U$ with no holomorphic logarithm, and relate this to $\pi_1(U)$.
\end{exercise}

\begin{exercise}\label{complex-analysis:exercise-mobius-disc-half-plane}
Find the M\"obius transformation sending $1$, $i$, $-1$ to $0$, $1$, $\infty$, and show that it maps $\mathbb{D}$ onto $\mathbb{H}$. Find a M\"obius transformation mapping the half-plane $\{\operatorname{Re}z > 0\}$ onto $\mathbb{D}$
with $1 \mapsto 0$.
\end{exercise}

\begin{solution}
By Theorem~\ref{complex-analysis:theorem-three-points}, $S(z) = \frac{(z - 1)(i + 1)}{(z + 1)(i - 1)}$, and $\frac{1 + i}{i - 1} = -i$, so $S(z) = i\frac{1 - z}{1 + z}$. It maps the unit circle, through $1$, $i$, $-1$, onto
$\RR \cup \{\infty\}$, so it maps $\mathbb{D}$ onto one of the two half-planes; as $S(0) = i$, it is $\mathbb{H}$. For the second part, $z \mapsto iz$ maps $\{\operatorname{Re}z > 0\}$ onto $\mathbb{H}$, and then the Cayley
transform gives $\frac{iz - i}{iz + i} = \frac{z - 1}{z + 1}$, which sends $1$ to $0$.
\end{solution}

\begin{exercise}\label{complex-analysis:exercise-cross-ratio-permutations}
Let $\lambda = (z, z_1, z_2, z_3)$. Show that permuting the four points changes the cross-ratio into one of $\lambda$, $1 - \lambda$, $1/\lambda$, $1/(1 - \lambda)$, $(\lambda - 1)/\lambda$, $\lambda/(\lambda - 1)$, and that
each value is obtained from four of the $24$ permutations.
\end{exercise}

\begin{solution}
We may assume the four points are finite (move them by a M\"obius map, Corollary~\ref{complex-analysis:corollary-cross-ratio}(1)). The permutations $(z\,z_1)(z_2\,z_3)$, $(z\,z_2)(z_1\,z_3)$ and $(z\,z_3)(z_1\,z_2)$ leave the
cross-ratio unchanged: for instance $(z_1, z, z_3, z_2) = \frac{(z_1 - z)(z_3 - z_2)}{(z_1 - z_2)(z_3 - z)}$, and numerator and denominator are those of $(z, z_1, z_2, z_3)$, each with two signs changed; the other two
are similar. With the identity they form a normal subgroup $V$ of order $4$, and every permutation is the product of an element of $V$ and a permutation fixing the first entry, so it suffices to permute
$z_1, z_2, z_3$. Moving them to $0, 1, \infty$, so that $z = \lambda$, and using the formulas of Theorem~\ref{complex-analysis:theorem-three-points} for points at $\infty$:
\[
(\lambda, 0, 1, \infty) = \lambda, \quad (\lambda, 1, 0, \infty) = 1 - \lambda, \quad (\lambda, 0, \infty, 1) = \frac{\lambda}{\lambda - 1}, \quad (\lambda, 1, \infty, 0) = \frac{\lambda - 1}{\lambda}, \quad
(\lambda, \infty, 0, 1) = \frac{1}{1 - \lambda}, \quad (\lambda, \infty, 1, 0) = \frac1\lambda.
\]
Each value is obtained from a coset of $V$, that is, from four permutations.
\end{solution}

\begin{exercise}\label{complex-analysis:exercise-mobius-classify}
Classify the following M\"obius transformations as parabolic, elliptic, hyperbolic or loxodromic, and find their fixed points: $z \mapsto \frac{z}{z + 1}$, $z \mapsto \frac{2z + 1}{z + 1}$, $z \mapsto 2iz$,
$z \mapsto -\frac{1}{z}$, $z \mapsto \frac{z - i}{z + i}$. Show that a M\"obius transformation of finite order $n \geq 2$ is elliptic and conjugate to $z \mapsto e^{2\pi ij/n}z$ for some $j$ prime to $n$.
\end{exercise}

\begin{solution}
Compute $t^2 = (\operatorname{tr}A)^2/\det A$, which is the invariant for the normalised matrix. $\frac{z}{z + 1}$: $t^2 = 4$, parabolic, fixed point $0$. $\frac{2z + 1}{z + 1}$: $t^2 = 9$, hyperbolic, fixed points
$z^2 - z - 1 = 0$, $z = (1 \pm \sqrt5)/2$. $2iz$: $t^2 = (1 + 2i)^2/(2i) = (-3 + 4i)/(2i) = 2 + \frac32i$, loxodromic, fixed points $0$, $\infty$. $-1/z$: $t^2 = 0$, elliptic, fixed points $\pm i$. $\frac{z - i}{z + i}$:
$t^2 = 4/(2i) = -2i$, loxodromic, fixed points the roots of $z^2 + (i - 1)z + i = 0$. If $M^n = \id$ with $M \neq \id$, then $M$ is not parabolic (the iterates of $z + 1$ are $z + n$), so it is conjugate to $z \mapsto kz$
with $k^n = 1$, $k \neq 1$, which is elliptic; replacing $n$ by the exact order, $k$ is a primitive $n$-th root of unity.
\end{solution}

\begin{exercise}\label{complex-analysis:exercise-joukowski}
Let $J(z) = \frac12(z + z^{-1})$. Show that $J$ is a rational function of degree $2$ with critical points $\pm1$, that it maps the circle $|z| = r \neq 1$ onto an ellipse with foci $\pm1$ and the unit circle onto
$[-1, 1]$, and that it maps $\mathbb{D} \setminus \{0\}$ biholomorphically onto $\CC \setminus [-1, 1]$. (This \emph{Joukowski map} was used to model the flow around aircraft wings.)
\end{exercise}

\begin{solution}
$J = (z^2 + 1)/(2z)$ has degree $2$, and $J^{\prime}(z) = \frac12(1 - z^{-2})$ vanishes at $\pm1$. For $z = re^{i\theta}$, $J(z) = \frac12(r + r^{-1})\cos\theta + \frac i2(r - r^{-1})\sin\theta$, an ellipse with semi-axes
$a = \frac12(r + r^{-1})$ and $b = \frac12|r - r^{-1}|$, so $a^2 - b^2 = 1$ and the foci are $\pm1$; for $r = 1$ it is $\cos\theta \in [-1, 1]$. Next $J(z) = J(w)$ if and only if $(z - w)(1 - \frac{1}{zw}) = 0$, that is,
$w = z$ or $w = 1/z$; so $J$ is injective on $\mathbb{D} \setminus \{0\}$. Every $w \in \CC$ is taken twice counted with multiplicity (Proposition~\ref{complex-analysis:proposition-rational-degree}), by $z$ and $1/z$;
if $w \notin [-1, 1]$, neither lies on the unit circle, so exactly one lies in $\mathbb{D} \setminus \{0\}$. The inverse is holomorphic, since $J^{\prime} \neq 0$ on $\mathbb{D} \setminus \{0\}$.
\end{solution}

\begin{exercise}\label{complex-analysis:exercise-apollonius}
Let $a \neq b$ in $\CC$ and $k > 0$. Show that $\{z : |z - a| = k|z - b|\}$ is a circle for $k \neq 1$ and the perpendicular bisector of $a$ and $b$ for $k = 1$ (a \emph{circle of Apollonius}). Deduce that the
images of the circles $|w| = k$ under the M\"obius map $w \mapsto \frac{a - bw}{1 - w}$, inverse of $z \mapsto \frac{z - a}{z - b}$, are these circles.
\end{exercise}

\begin{solution}
Squaring, $|z|^2 - 2\operatorname{Re}(z\bar a) + |a|^2 = k^2(|z|^2 - 2\operatorname{Re}(z\bar b) + |b|^2)$, that is,
$(1 - k^2)|z|^2 - 2\operatorname{Re}\bigl(z(\bar a - k^2\bar b)\bigr) + |a|^2 - k^2|b|^2 = 0$, which is of the form~\eqref{complex-analysis:equation-generalised-circle}; it is a line, the bisector, if $k = 1$. For $k \neq 1$ the
set is nonempty and not a point (it meets the line through $a$ and $b$ in two points, dividing the segment internally and externally in the ratio $k$), so it is a circle. The second statement is the first, since
$|w| = k$ means $|z - a| = k|z - b|$ for $z$ the preimage of $w$.
\end{solution}

\begin{exercise}\label{complex-analysis:exercise-residue-infinity}
Compute $\int_{|z| = 3}\frac{z^{5}\,dz}{(z - 1)(z^5 + 2)}$ using the residue at $\infty$.
\end{exercise}

\begin{solution}
All poles lie in $|z| < 3$. For large $|z|$, $\frac{z^5}{(z - 1)(z^5 + 2)} = \frac{1}{z - 1}\cdot\frac{1}{1 + 2z^{-5}} = z^{-1}(1 + z^{-1} + \cdots)(1 - 2z^{-5} + \cdots)$, whose coefficient of $z^{-1}$ is $1$. So
$\operatorname{Res}_\infty = -1$ and, by Proposition~\ref{complex-analysis:proposition-residues-sphere}, the integral is $2\pi i\sum_{a \in \CC}\operatorname{Res}_a = -2\pi i\operatorname{Res}_\infty = 2\pi i$.
\end{solution}

\begin{exercise}\label{complex-analysis:exercise-constant-modulus}
Let $f = u + iv$ be holomorphic on a connected open set $U$. Show that $f$ is constant if any of the following holds: (1) $u$ is constant; (2) $f$ takes only real values; (3) $|f|$ is constant; (4) $\bar f$ is
also holomorphic.
\end{exercise}

\begin{solution}
If $f^{\prime} = 0$ on $U$, then $f$ is constant on every disc in $U$ (Lemma~\ref{complex-analysis:lemma-path-integral}, integrating along segments), so the set where $f$ takes a given value is open, as is its
complement, and $f$ is constant since $U$ is connected. (1) If $u$ is constant, $\partial_xu = \partial_yu = 0$, so $\partial_xv = -\partial_yu = 0$ by the Cauchy--Riemann equations, and
$f^{\prime} = \partial_xu + i\partial_xv = 0$. (2) Then $v = 0$, and $-if = v - iu$ has constant real part; apply (1). (3) If $|f|^2$ is constant, then $0 = \Delta|f|^2 = 4|f^{\prime}|^2$ by
Example~\ref{complex-analysis:example-harmonic}(3). (4) Then $f + \bar f = 2u$ is holomorphic and real valued; by (2) $u$ is constant, and (1) applies.
\end{solution}

\begin{exercise}\label{complex-analysis:exercise-polar-cauchy-riemann}
(\emph{Cauchy--Riemann equations in polar coordinates.}) Let $f = u + iv$ be real differentiable on an open set $U \subseteq \CC^\times$, and write $z = re^{i\theta}$. Show that $f$ is holomorphic if and
only if
\[
\partial_ru = \frac1r\partial_\theta v, \qquad \partial_rv = -\frac1r\partial_\theta u,
\]
and that then $f^{\prime}(z) = e^{-i\theta}\partial_rf$. Deduce that the principal logarithm $\operatorname{Log}z = \log r + i\theta$, $-\pi < \theta < \pi$, is holomorphic on $\CC \setminus (-\infty, 0]$ with
derivative $1/z$.
\end{exercise}

\begin{solution}
By the chain rule, $\partial_r = \cos\theta\,\partial_x + \sin\theta\,\partial_y$ and $\partial_\theta = -r\sin\theta\,\partial_x + r\cos\theta\,\partial_y$ on functions of $(x, y) = (r\cos\theta, r\sin\theta)$. Hence
\[
\partial_\theta f - ir\,\partial_rf = r\bigl((-\sin\theta - i\cos\theta)\partial_xf + (\cos\theta - i\sin\theta)\partial_yf\bigr) = -ire^{-i\theta}\bigl(\partial_xf + i\,\partial_yf\bigr).
\]
The Cauchy--Riemann equations say $\partial_xf + i\,\partial_yf = 0$ (its real and imaginary parts are $\partial_xu - \partial_yv$ and $\partial_xv + \partial_yu$). So $f$ is holomorphic if and only if
$\partial_\theta f = ir\,\partial_rf$, whose real and imaginary parts are the stated equations. Then $\partial_rf = \cos\theta\,f^{\prime} + \sin\theta\,if^{\prime} = e^{i\theta}f^{\prime}$. For
$\operatorname{Log}$, which is real differentiable because $(r, \theta)$ are smooth functions of $(x, y)$ on the slit plane, $u = \log r$ and $v = \theta$ give $\partial_ru = \frac1r = \frac1r\partial_\theta v$ and
$\partial_rv = 0 = -\frac1r\partial_\theta u$, and $\operatorname{Log}^{\prime}(z) = e^{-i\theta}/r = 1/z$.
\end{solution}

\begin{exercise}\label{complex-analysis:exercise-circle-monomials}
For integers $m, n \geq 0$ compute $\int_\gamma z^n\bar z^m\,dz$ over the unit circle $\gamma(t) = e^{2\pi it}$. Deduce that a polynomial $p(z, \bar z) = \sum c_{nm}z^n\bar z^m$ satisfies
$\int_\gamma p\,dz = 2\pi i\sum_nc_{n,n+1}$; in particular this integral can be nonzero only if $p$ really involves $\bar z$.
\end{exercise}

\begin{solution}
On the circle $\bar z = 1/z$, so the integrand is $z^{n-m}$, and by Lemma~\ref{complex-analysis:lemma-path-integral} the integral is $2\pi i$ if $n - m = -1$ and $0$ otherwise. The formula for $p$ follows by
linearity, and if $p$ does not involve $\bar z$ then $c_{n,n+1} = 0$ for all $n$.
\end{solution}

\begin{exercise}\label{complex-analysis:exercise-cauchy-kernel-dbar}
Show that $1/z$ is locally integrable on $\CC$, and that for every smooth $\varphi \colon \CC \to \CC$ with compact support
\[
\int_\CC\frac1z\,\frac{\partial\varphi}{\partial\bar z}\,dA = -\pi\varphi(0).
\]
So $1/z$ does not satisfy the weak equation of Proposition~\ref{complex-analysis:proposition-weak-cauchy-riemann} on $\CC$, although it is holomorphic on $\CC^\times$: the weak equation sees the pole. In the language of
distributions, $\frac{\partial}{\partial\bar z}\frac{1}{\pi z} = \delta_0$, and $\frac1{\pi z}$ is a fundamental solution of $\partial/\partial\bar z$. (Hint: in polar coordinates
$\frac{\partial}{\partial\bar z} = \frac{e^{i\theta}}{2}\bigl(\partial_r + \frac ir\partial_\theta\bigr)$.)
\end{exercise}

\begin{solution}
In polar coordinates $z = re^{i\theta}$ the area element is $dA = r\,dr\,d\theta$ (Corollary~\ref{measure-theory:corollary-polar} for $n = 2$, the measure $\mathcal{H}^1$ on the unit circle being $d\theta$), so
$\int_{|z| < R}|z|^{-1}\,dA = \int_0^{2\pi}\int_0^R dr\,d\theta = 2\pi R < \infty$: $1/z$ is locally integrable. By the chain rule, $\partial_r = \cos\theta\,\partial_x + \sin\theta\,\partial_y$ and
$\partial_\theta = -r\sin\theta\,\partial_x + r\cos\theta\,\partial_y$, so $\partial_r + \frac ir\partial_\theta = e^{-i\theta}\partial_x + ie^{-i\theta}\partial_y = 2e^{-i\theta}\frac{\partial}{\partial\bar z}$, which is the hint.
Hence $\frac1z\frac{\partial\varphi}{\partial\bar z} = \frac{1}{2r}\bigl(\partial_r\varphi + \frac ir\partial_\theta\varphi\bigr)$, and by dominated convergence (the integrand is integrable) and Fubini's theorem,
\[
\int_\CC\frac1z\frac{\partial\varphi}{\partial\bar z}\,dA = \lim_{\varepsilon \to 0}\frac12\int_\varepsilon^\infty\int_0^{2\pi}\Bigl(\partial_r\varphi + \frac ir\partial_\theta\varphi\Bigr)d\theta\,dr .
\]
The $\theta$-integral of $\partial_\theta\varphi$ vanishes, $\varphi(re^{i\theta})$ being $2\pi$-periodic in $\theta$, and $\int_\varepsilon^\infty\partial_r\varphi\,dr = -\varphi(\varepsilon e^{i\theta})$ since $\varphi$ has compact support.
So the integral is $\lim_{\varepsilon \to 0}\bigl(-\frac12\int_0^{2\pi}\varphi(\varepsilon e^{i\theta})\,d\theta\bigr) = -\pi\varphi(0)$, by continuity of $\varphi$.
\end{solution}
